% =====================================================================
%  utpnotes.tex
%  Shared preamble for the UTP teaching-assistant course notes.
%
%  All eight books include this file. It owns everything they have in
%  common, so a change here reaches every course at once. Anything a
%  single book needs differently is overridden in that book's own
%  main-<code>.tex, AFTER the \input line.
%
%  ------------------------------------------------------------------
%  HOW A COURSE USES THIS FILE
%  ------------------------------------------------------------------
%      \documentclass[11pt,oneside,a4paper]{report}
%
%      \def\utprunninghead{Machine Learning Notes}   % see section 1
%      \def\utplabtotal{10}
%
%      % =====================================================================
%  utpnotes.tex
%  Shared preamble for the UTP teaching-assistant course notes.
%
%  All eight books include this file. It owns everything they have in
%  common, so a change here reaches every course at once. Anything a
%  single book needs differently is overridden in that book's own
%  main-<code>.tex, AFTER the \input line.
%
%  ------------------------------------------------------------------
%  HOW A COURSE USES THIS FILE
%  ------------------------------------------------------------------
%      \documentclass[11pt,oneside,a4paper]{report}
%
%      \def\utprunninghead{Machine Learning Notes}   % see section 1
%      \def\utplabtotal{10}
%
%      % =====================================================================
%  utpnotes.tex
%  Shared preamble for the UTP teaching-assistant course notes.
%
%  All eight books include this file. It owns everything they have in
%  common, so a change here reaches every course at once. Anything a
%  single book needs differently is overridden in that book's own
%  main-<code>.tex, AFTER the \input line.
%
%  ------------------------------------------------------------------
%  HOW A COURSE USES THIS FILE
%  ------------------------------------------------------------------
%      \documentclass[11pt,oneside,a4paper]{report}
%
%      \def\utprunninghead{Machine Learning Notes}   % see section 1
%      \def\utplabtotal{10}
%
%      % =====================================================================
%  utpnotes.tex
%  Shared preamble for the UTP teaching-assistant course notes.
%
%  All eight books include this file. It owns everything they have in
%  common, so a change here reaches every course at once. Anything a
%  single book needs differently is overridden in that book's own
%  main-<code>.tex, AFTER the \input line.
%
%  ------------------------------------------------------------------
%  HOW A COURSE USES THIS FILE
%  ------------------------------------------------------------------
%      \documentclass[11pt,oneside,a4paper]{report}
%
%      \def\utprunninghead{Machine Learning Notes}   % see section 1
%      \def\utplabtotal{10}
%
%      \input{../.shared/utpnotes}
%      \graphicspath{{../.shared/}}
%
%  Configuration goes BEFORE the \input; overrides go AFTER it.
%  Paths are relative, so compile from inside the course folder.
%
%  Adding a new course? Copy .shared/TEMPLATE-main.tex and follow
%  .shared/ADDING-A-COURSE.md. Every step is a numbered checklist.
%
%  ------------------------------------------------------------------
%  CONTENTS
%  ------------------------------------------------------------------
%    1  Course configuration     the knobs a book may set
%    2  Packages                 the shared package set
%    3  Colour palette           one palette, plus per-course aliases
%    4  Page layout and tables   margins, lists, table styling
%    5  Table of contents        entry formatting for parts/chapters/labs
%    6  Headings                 chapter badge, part pages, sections
%    7  Page furniture           running heads and footers
%    8  Learning boxes           key idea, pitfall, try it, and friends
%    9  Small commands           \term, \cmd, \chapterguide, \mlsteps
%   10  Lab milestone banner     the full-width "Lab N of M" page break
%   11  Stage divider            the part-opening page in longer books
%   12  Listing styles           the base code and terminal styles
% =====================================================================

% --------------------------------------------------------------------
%  1.  COURSE CONFIGURATION
% --------------------------------------------------------------------
%  The three knobs below, plus \utpstyleparts in section 5 and the box
%  titles in section 8, are the whole configuration surface. A book
%  sets what it needs BEFORE the \input and inherits the rest.
%
%    \utprunninghead   the left running head on every page.
%    \utpchapterlabel  the word in the chapter badge. CHAPTER by
%                      default; a book whose chapters are labs sets LAB.
%    \utplabtotal      the N in the "Lab 3 of N" banners.
%
%  \providecommand means "define unless the book already did", which is
%  what makes the \def-before-\input pattern work.

\providecommand{\utprunninghead}{Course Notes}
\providecommand{\utpchapterlabel}{CHAPTER}
\providecommand{\utplabtotal}{10}

% --------------------------------------------------------------------
%  2.  PACKAGES
% --------------------------------------------------------------------
%  The union of what the eight books need. A course that needs
%  something extra loads it in its own file, under the heading
%  "packages only this book needs" (algorithm2e, caption, float and
%  multirow are the only current examples).
%
%  The tikz library list is the union across all books, so a diagram
%  copied from one course compiles in another.

\usepackage[T1]{fontenc}
\IfFileExists{glyphtounicode.tex}{\input{glyphtounicode}\pdfgentounicode=1}{}
\usepackage[american]{babel}
\IfFileExists{sourcesanspro.sty}
  {\usepackage[default]{sourcesanspro}}
  {\usepackage[scaled]{helvet}\renewcommand{\familydefault}{\sfdefault}}
\DeclareFontShape{T1}{cmss}{b}{n}{<->ssub*cmss/bx/n}{}
\DeclareFontShape{TS1}{cmss}{b}{n}{<->ssub*cmss/bx/n}{}
\usepackage[margin=21mm,headheight=22pt]{geometry}
\usepackage{amsmath,amssymb,mathtools,bm}
\usepackage{microtype}
\usepackage{enumitem}
\usepackage{booktabs}
\usepackage{tabularx}
\usepackage{array}
\usepackage{ragged2e}
\usepackage{etoolbox}
\usepackage{titlesec}
\usepackage{tocloft}
\usepackage{fancyhdr}
\usepackage[table]{xcolor}
\usepackage{tcolorbox}
\tcbuselibrary{breakable,skins}
\usepackage[hidelinks]{hyperref}
\usepackage{bookmark}
\usepackage{listings}
\usepackage{graphicx}
\usepackage{tikz}
\usetikzlibrary{angles,arrows.meta,backgrounds,calc,chains,decorations.pathmorphing,decorations.pathreplacing,fit,matrix,patterns,positioning,quotes,shadows.blur,shapes.geometric,shapes.misc}
\usepackage{pgfplots}

% --------------------------------------------------------------------
%  3.  COLOUR PALETTE
% --------------------------------------------------------------------
%  One palette, defined once under canonical UTP* names. Use these
%  in any new material.
%
%  The existing chapters were written against course-prefixed names
%  (MLNavy, ADSPaleBlue, ...) in over two thousand places, so rather
%  than rewrite them the block below aliases each prefix onto the
%  canonical colours. All four sets are always defined, so a chapter
%  moved between courses still resolves.

\definecolor{UTPNavy}{HTML}{18324A}
\definecolor{UTPBlue}{HTML}{246B9E}
\definecolor{UTPTeal}{HTML}{138A80}
\definecolor{UTPGreen}{HTML}{2D7D46}
\definecolor{UTPOrange}{HTML}{C86B24}
\definecolor{UTPRed}{HTML}{B34843}
\definecolor{UTPGold}{HTML}{D4A13E}
\definecolor{UTPInk}{HTML}{263238}
\definecolor{UTPMuted}{HTML}{60717C}
\definecolor{UTPPaleBlue}{HTML}{EEF6FB}
\definecolor{UTPPaleTeal}{HTML}{EEF8F6}
\definecolor{UTPPaleOrange}{HTML}{FFF5EA}
\definecolor{UTPPaleRed}{HTML}{FFF0EF}
\definecolor{UTPPaleGray}{HTML}{F4F6F7}

% Course files were written against prefixed colour names. Rather than edit
% 2000+ references across the chapters, alias each prefix onto the canonical
% palette. Only these four prefixes appear anywhere in the books.
\def\utpmkalias#1{%
  \colorlet{#1Navy}{UTPNavy}\colorlet{#1Blue}{UTPBlue}%
  \colorlet{#1Teal}{UTPTeal}\colorlet{#1Green}{UTPGreen}%
  \colorlet{#1Orange}{UTPOrange}\colorlet{#1Red}{UTPRed}%
  \colorlet{#1Ink}{UTPInk}\colorlet{#1Muted}{UTPMuted}%
  \colorlet{#1PaleBlue}{UTPPaleBlue}\colorlet{#1PaleTeal}{UTPPaleTeal}%
  \colorlet{#1PaleOrange}{UTPPaleOrange}\colorlet{#1PaleRed}{UTPPaleRed}%
  \colorlet{#1PaleGray}{UTPPaleGray}%
}
\utpmkalias{ML}
\utpmkalias{ADS}
\utpmkalias{DS}
\utpmkalias{ERP}
\color{UTPInk}

% --------------------------------------------------------------------
%  4.  PAGE LAYOUT AND TABLES
% --------------------------------------------------------------------
%  Paragraph spacing, list indents, and the table look: alternating
%  row tint, top-aligned paragraph columns, ragged-right cells.
%
%  A book typeset tighter overrides \arraystretch, \tabcolsep or
%  \emergencystretch after the input. Data Communication is the one
%  that currently does.

\setlength{\parindent}{0pt}
\setlength{\parskip}{0.55em}
\setlist[itemize]{leftmargin=1.45em,itemsep=0.22em,topsep=0.25em}
\setlist[enumerate]{leftmargin=1.75em,itemsep=0.28em,topsep=0.3em}
\setlist[description]{font=\bfseries\color{UTPNavy},itemsep=0.25em,topsep=0.3em}
\renewcommand{\arraystretch}{1.34}
\setlength{\tabcolsep}{6.5pt}
\arrayrulecolor{UTPBlue!45}
% Keep paragraph-style columns top-aligned and ragged. Tables mix fixed p
% columns with flexible X columns; giving X the same baseline stops one
% column floating halfway down a multi-line row.
\renewcommand{\tabularxcolumn}[1]{p{#1}}
\makeatletter
\g@addto@macro\@arrayparboxrestore{\RaggedRight\arraybackslash}
\makeatother
\AtBeginEnvironment{tabular}{\small\setlength{\extrarowheight}{0.6pt}\rowcolors{2}{UTPPaleBlue!72}{white}}
\AtBeginEnvironment{tabularx}{\small\setlength{\extrarowheight}{0.6pt}\rowcolors{2}{UTPPaleBlue!72}{white}}
\AtBeginEnvironment{figure}{\centering}
\emergencystretch=1.5em
\hbadness=10000
\hfuzz=8pt

\pgfplotsset{
  compat=1.18,
  every axis/.append style={
    axis line style={draw=UTPMuted!85,line width=0.55pt},
    tick style={draw=UTPMuted!85},
    tick label style={font=\scriptsize\color{UTPMuted}},
    label style={font=\small\color{UTPNavy}},
    title style={font=\small\bfseries\color{UTPNavy}},
    legend style={font=\scriptsize,draw=UTPBlue!22,fill=white,rounded corners=1pt}
  },
  every axis plot/.append style={line width=1.15pt}
}
\tikzset{every picture/.append style={line cap=round,line join=round}}

% --------------------------------------------------------------------
%  5.  TABLE OF CONTENTS
% --------------------------------------------------------------------
%  Formatting for part, chapter and lab entries. The lab level is a
%  custom tocloft list so lab banners appear in the contents.
%
%  Part styling is skipped when \utpstyleparts is 0, which is what a
%  book without \part divisions should set.

\addto\captionsamerican{\renewcommand{\contentsname}{Learning Path}}
\newlistentry{lab}{toc}{-1}
\providecommand{\utpstyleparts}{1}
\ifnum\utpstyleparts=1
\renewcommand{\cftpartpresnum}{}
\renewcommand{\cftpartaftersnum}{}
\setlength{\cftpartnumwidth}{5.2em}
\setlength{\cftbeforepartskip}{0.38em}
\renewcommand{\cftpartfont}{\normalsize\bfseries\color{UTPNavy}}
\renewcommand{\cftpartpagefont}{\normalsize\bfseries\color{UTPNavy}}
\renewcommand{\cftpartleader}{\hfill}
\fi
\renewcommand{\cftchappresnum}{Chapter~}
\renewcommand{\cftchapaftersnum}{\quad}
\setlength{\cftchapindent}{1.5em}
\setlength{\cftchapnumwidth}{6em}
\setlength{\cftbeforechapskip}{0pt}
\renewcommand{\cftchapfont}{\normalfont}
\renewcommand{\cftchappagefont}{\normalfont}
\renewcommand{\cftchapleader}{\cftdotfill{\cftdotsep}}
\setlength{\cftlabindent}{3em}
\setlength{\cftlabnumwidth}{0em}
\setlength{\cftbeforelabskip}{0.10em}
\renewcommand{\cftlabfont}{\bfseries\color{UTPOrange!90!black}}
\renewcommand{\cftlabpagefont}{\bfseries\color{UTPOrange!90!black}}
\renewcommand{\cftlableader}{\color{UTPOrange!65!UTPMuted}\cftdotfill{\cftlabdotsep}}

% --------------------------------------------------------------------
%  6.  HEADINGS
% --------------------------------------------------------------------
%  The chapter badge (a coloured block with a word and the number),
%  part pages, and section spacing.
%
%  The word in the badge is \utpchapterlabel; Data Communication sets
%  it to LAB because each of its chapters is a lab.

\titleformat{\chapter}[display]
  {\bfseries\color{UTPNavy}}
  {\filleft\colorbox{UTPBlue}{\parbox{2.55cm}{\centering\color{white}\Large
    \utpchapterlabel\\[-0.15em]\Huge\thechapter}}}
  {0.65em}{\Huge\raggedright}
  [\vspace{0.25em}{\color{UTPBlue}\titlerule[1.2pt]}]
\renewcommand{\thepart}{Stage~\arabic{part}}
\titleformat{\part}[display]
  {\centering\color{UTPNavy}\bfseries}
  {\Large\color{UTPTeal}\thepart}{0.45em}{\Huge}
  [\vspace{0.55em}\color{UTPTeal}\rule{0.45\linewidth}{1.5pt}]
\titleformat{\section}{\Large\bfseries\color{UTPNavy}}{\thesection}{0.55em}{}
\titleformat{\subsection}{\large\bfseries\color{UTPBlue}}{\thesubsection}{0.5em}{}
\titlespacing*{\chapter}{0pt}{-1.8em}{1.0em}
\titlespacing*{\part}{0pt}{0pt}{1em}
\titlespacing*{\section}{0pt}{1.2em}{0.35em}
\titlespacing*{\subsection}{0pt}{0.95em}{0.2em}

% --------------------------------------------------------------------
%  7.  PAGE FURNITURE
% --------------------------------------------------------------------
%  Running heads and footers. The left header is \utprunninghead;
%  the right one follows the current chapter.

\pagestyle{fancy}
\fancyhf{}
\lhead{\small\color{UTPMuted}\utprunninghead}
\rhead{\small\color{UTPMuted}\nouppercase{\leftmark}}
\cfoot{\small\color{UTPMuted}\thepage}
\renewcommand{\headrulewidth}{0.5pt}
\renewcommand{\headrule}{\hbox to\headwidth{\color{UTPPaleBlue}\leaders\hrule height \headrulewidth\hfill}}
\renewcommand{\chaptermark}[1]{\markboth{#1}{}}
\fancypagestyle{plain}{%
  \fancyhf{}%
  \cfoot{\small\color{UTPMuted}\thepage}%
  \renewcommand{\headrulewidth}{0pt}%
}

% --------------------------------------------------------------------
%  8.  LEARNING BOXES
% --------------------------------------------------------------------
%  The coloured callouts the chapters use. Every title is a macro, so
%  a book can rename a box without restating its geometry:
%
%      \def\utptitlepitfall{Technical note}   % before the input
%
%  To change how a box LOOKS in one book, use \renewtcolorbox after
%  the input instead. To add a box only one book needs, declare it
%  there with \newtcolorbox.

% Box titles -- override any of these before the input line.
\providecommand{\utptitlekeyidea}{Key idea}
\providecommand{\utptitleworked}{Worked example}
\providecommand{\utptitlepitfall}{Common mistake}
\providecommand{\utptitletryit}{Try it yourself}
\providecommand{\utptitlequickcheck}{Check your understanding}
\providecommand{\utptitlerecap}{Chapter recap}
\providecommand{\utptitledeliverable}{Lab deliverable}
\providecommand{\utptitlelabmode}{How much is scaffolded?}
\providecommand{\utptitleintuition}{Plain idea}
\providecommand{\utptitlefirstread}{Core path}
\providecommand{\utptitleglance}{Chapter at a glance}
\providecommand{\utptitledebug}{Debugging task}

% Explaining an idea
\newtcolorbox{keyidea}{enhanced,breakable,colback=UTPPaleTeal,colframe=UTPTeal, boxrule=0.8pt,arc=2mm,left=2.5mm,right=2.5mm,top=1.8mm,bottom=1.8mm, fonttitle=\bfseries\color{UTPNavy},title=\utptitlekeyidea}
\newtcolorbox{workedexample}{enhanced,breakable,colback=UTPPaleBlue,colframe=UTPBlue, boxrule=0.8pt,arc=2mm,left=2.5mm,right=2.5mm,top=1.8mm,bottom=1.8mm, fonttitle=\bfseries\color{UTPNavy},title=\utptitleworked}
\newtcolorbox{pitfall}{enhanced,breakable,colback=UTPPaleRed,colframe=UTPRed, boxrule=0.8pt,arc=2mm,left=2.5mm,right=2.5mm,top=1.8mm,bottom=1.8mm, fonttitle=\bfseries\color{UTPNavy},title=\utptitlepitfall}
% Asking the reader to do something
\newtcolorbox{tryit}{enhanced,breakable,colback=white,colframe=UTPTeal, boxrule=0.9pt,arc=2mm,left=2.5mm,right=2.5mm,top=1.8mm,bottom=1.8mm, fonttitle=\bfseries\color{white},colbacktitle=UTPTeal,title=\utptitletryit}
\newtcolorbox{quickcheck}{enhanced,breakable,colback=white,colframe=UTPOrange, boxrule=0.9pt,arc=2mm,left=2.5mm,right=2.5mm,top=1.8mm,bottom=1.8mm, fonttitle=\bfseries\color{white},colbacktitle=UTPOrange,title=\utptitlequickcheck}
\newtcolorbox{chapterrecap}{enhanced,breakable,colback=UTPPaleTeal,colframe=UTPGreen, boxrule=0.9pt,arc=2mm,left=2.5mm,right=2.5mm,top=1.8mm,bottom=1.8mm, fonttitle=\bfseries\color{white},colbacktitle=UTPGreen,title=\utptitlerecap}
% Lab structure
\newtcolorbox{labdeliverable}{enhanced,breakable,colback=white,colframe=UTPGreen, boxrule=0.9pt,arc=2mm,left=2.6mm,right=2.6mm,top=1.9mm,bottom=1.9mm, fonttitle=\bfseries\color{white},colbacktitle=UTPGreen,title=\utptitledeliverable}
\newtcolorbox{labmode}{enhanced,breakable,colback=UTPPaleBlue!52,colframe=UTPBlue, boxrule=0.75pt,arc=2mm,left=2.6mm,right=2.6mm,top=1.7mm,bottom=1.7mm, fonttitle=\bfseries\color{white},colbacktitle=UTPBlue,title=\utptitlelabmode}
% Softer, lighter-weight asides
\newtcolorbox{intuition}{enhanced,breakable,colback=white,colframe=UTPOrange!42, boxrule=0.5pt,arc=2mm,left=2.5mm,right=2.5mm,top=1.5mm,bottom=1.5mm, fonttitle=\bfseries\color{UTPOrange!85!black},title=\utptitleintuition}
\newtcolorbox{firstread}{enhanced,breakable,colback=white,colframe=UTPBlue!42, boxrule=0.5pt,arc=2mm,left=2.5mm,right=2.5mm,top=1.5mm,bottom=1.5mm, fonttitle=\bfseries\color{UTPBlue}, title=\utptitlefirstread}
\newtcolorbox{debugtask}{enhanced,breakable,colback=UTPPaleRed!58,colframe=UTPRed, boxrule=0.85pt,arc=2mm,left=2.6mm,right=2.6mm,top=1.9mm,bottom=1.9mm, fonttitle=\bfseries\color{white},colbacktitle=UTPRed,title=\utptitledebug}
\newtcolorbox{analogy}{enhanced,breakable,colback=white,colframe=UTPOrange!38, boxrule=0.5pt,arc=2mm,left=2.5mm,right=2.5mm,top=1.5mm,bottom=1.5mm, fonttitle=\bfseries\color{UTPOrange!85!black},title=Analogy}
\newtcolorbox{coursechoice}{enhanced,breakable,colback=white,colframe=UTPNavy, boxrule=0.8pt,arc=2mm,left=2.8mm,right=2.8mm,top=2mm,bottom=2mm, fonttitle=\bfseries\color{white},colbacktitle=UTPNavy,title=Course choice}
\newtcolorbox{decisionmemo}{enhanced,breakable,colback=UTPPaleOrange!55,colframe=UTPOrange, boxrule=0.85pt,arc=2mm,left=2.6mm,right=2.6mm,top=1.9mm,bottom=1.9mm, fonttitle=\bfseries\color{white},colbacktitle=UTPOrange,title=Decision memo}
\newtcolorbox{libnote}{enhanced,breakable,colback=white,colframe=UTPMuted!48, boxrule=0.5pt,arc=2mm,left=2.5mm,right=2.5mm,top=1.5mm,bottom=1.5mm, fonttitle=\bfseries\color{UTPNavy},title=Reference note}
\newtcolorbox{miniproject}{enhanced,breakable,colback=UTPPaleTeal!52,colframe=UTPGreen, boxrule=0.9pt,arc=2mm,left=2.6mm,right=2.6mm,top=1.9mm,bottom=1.9mm, fonttitle=\bfseries\color{white},colbacktitle=UTPGreen,title=Mini-project}
% Boxes that take a title argument
\newtcolorbox{modelcard}[1]{enhanced,breakable,colback=white,colframe=UTPBlue, boxrule=0.9pt,arc=2mm,left=2.5mm,right=2.5mm,top=1.8mm,bottom=1.8mm, fonttitle=\bfseries\color{white},colbacktitle=UTPBlue,title=#1}
\newtcolorbox{optionaltopic}{enhanced,breakable,colback=white,colframe=UTPMuted!58, boxrule=0.5pt,arc=2mm,left=2.5mm,right=2.5mm,top=1.5mm,bottom=1.5mm, fonttitle=\bfseries\color{UTPMuted},title=Beyond the core}
\newtcolorbox{secondpass}{enhanced,breakable,colback=UTPPaleGray!82,colframe=UTPMuted, boxrule=0.75pt,arc=2mm,left=2.6mm,right=2.6mm,top=1.9mm,bottom=1.9mm, fonttitle=\bfseries\color{white},colbacktitle=UTPMuted,title=Second pass}
\newtcolorbox{labproject}[1]{enhanced,breakable,colback=white,colframe=UTPBlue, boxrule=0.9pt,arc=2mm,left=2.6mm,right=2.6mm,top=1.9mm,bottom=1.9mm, fonttitle=\bfseries\color{white},colbacktitle=UTPBlue,title={#1}}
% Frames for a figure or a table
\newtcolorbox{mlfigure}{enhanced,width=\linewidth,colback=white,colframe=UTPBlue!28, boxrule=0.7pt,arc=2mm, left=2.2mm,right=2.2mm,top=2.2mm,bottom=1.8mm, before skip=0.75em,after skip=0.85em, before upper={\centering}, halign=center}
\newtcolorbox{mltable}{enhanced,width=\linewidth,colback=white,colframe=UTPTeal!32, boxrule=0.7pt,arc=2mm, left=2.2mm,right=2.2mm,top=1.8mm,bottom=1.8mm, before skip=0.75em,after skip=0.85em, halign=center}

% --------------------------------------------------------------------
%  9.  SMALL COMMANDS
% --------------------------------------------------------------------
%  \term for a defined term, \cmd for inline code, \chapref and
%  \secref for cross-references, and the mlsteps numbered-step list.
%
%      \term{overfitting}          a term being defined
%      \cmd{git status}            inline code
%      \chapref{ch:trees}          "Chapter 12"
%      \secref{sec:pruning}        "Section 12.3"
%      \chapterguide{prior}{outcomes}   the box that opens a chapter
%      \begin{mlsteps} ... \end{mlsteps}  a numbered Step 1/Step 2 list

\newcommand{\term}[1]{\textbf{\color{UTPNavy}#1}}
\newcommand{\cmd}[1]{\texttt{#1}}
\newcommand{\chapref}[1]{Chapter~\ref{#1}}
\newcommand{\secref}[1]{Section~\ref{#1}}
\newcommand{\chapterguide}[2]{%
  \begin{tcolorbox}[enhanced,breakable,colback=white,colframe=UTPNavy,boxrule=0.8pt,
    arc=2mm,left=2.5mm,right=2.5mm,top=2mm,bottom=2mm,
    title=\utptitleglance,fonttitle=\bfseries\color{white},colbacktitle=UTPNavy]
    \textbf{Before you start:} #1\par
    \textbf{By the end, you can:} #2
  \end{tcolorbox}}
\newenvironment{mlsteps}
  {\begin{enumerate}[label=\textcolor{UTPBlue}{\bfseries Step~\arabic*:},leftmargin=4.5em,itemsep=0.38em]}
  {\end{enumerate}}

% --------------------------------------------------------------------
%  10.  LAB MILESTONE BANNER
% --------------------------------------------------------------------
%  The full-width page break that opens each lab. Called from the lab
%  files as \labsection{number}{title}{label}.
%
%  The "of N" comes from \utplabtotal, so a book states its lab count
%  once in its configuration rather than in every lab file.

\newcommand{\labsection}[3]{%
  \clearpage
  \phantomsection
  \addcontentsline{toc}{lab}{Lab #1 of \utplabtotal\space--- #2}%
  \markboth{Lab #1 of \utplabtotal\space--- #2}{}%
  \begin{tcolorbox}[enhanced,width=\linewidth,colback=UTPNavy,colframe=UTPNavy,
    boxrule=0pt,arc=2mm,left=5mm,right=5mm,top=5mm,bottom=5mm,
    before skip=0pt,after skip=1.2em]
    {\large\bfseries\color{UTPTeal!35!white} PRACTICAL MILESTONE\par}
    \vspace{0.25em}
    {\Huge\bfseries\color{white} Lab #1 of \utplabtotal\par}
    \vspace{0.20em}
    {\Large\bfseries\color{white} #2\par}
  \end{tcolorbox}
  \label{#3}%
}

% --------------------------------------------------------------------
%  11.  STAGE DIVIDER
% --------------------------------------------------------------------
%  A full-page divider used by the longer books to open a part:
%  \substage{3A}{Supervised Workflow}{one-line summary}.
%  Harmless to leave unused.

\newcommand{\substage}[3]{%
  \cleardoublepage
  \thispagestyle{empty}
  \phantomsection
  \addcontentsline{toc}{chapter}{Stage #1 --- #2}
  \begin{center}
    \vspace*{0.20\textheight}
    {\Large\bfseries\color{UTPTeal} Stage #1\par}
    \vspace{0.45em}
    {\Huge\bfseries\color{UTPNavy} #2\par}
    \vspace{0.65em}
    {\color{UTPTeal}\rule{0.46\linewidth}{1.5pt}\par}
    \vspace{1.1em}
    \begin{minipage}{0.76\linewidth}
      \centering\large\color{UTPMuted} #3
    \end{minipage}
  \end{center}
  \cleardoublepage
}

% --------------------------------------------------------------------
%  12.  LISTING STYLES
% --------------------------------------------------------------------
%  Two base styles. A course builds its own on top of them so the
%  colours and frame stay consistent:
%
%      \lstdefinestyle{mlpython}{style=utpcode,language=Python}
%      \lstdefinestyle{mlterminal}{style=utpterminal}

\lstdefinestyle{utpcode}{
  basicstyle=\ttfamily\scriptsize,
  keywordstyle=\bfseries\color{UTPBlue},
  commentstyle=\itshape\color{UTPMuted},
  stringstyle=\color{UTPGreen!70!black},
  backgroundcolor=\color{UTPPaleGray!72},
  frame=single,
  rulecolor=\color{UTPBlue!28},
  framesep=5pt,
  breaklines=true,
  breakatwhitespace=false,
  showstringspaces=false,
  keepspaces=true,
  columns=fullflexible,
  upquote=true,
  tabsize=4,
  aboveskip=0.7em,
  belowskip=0.8em
}
\lstdefinestyle{utpterminal}{
  basicstyle=\ttfamily\scriptsize,
  backgroundcolor=\color{UTPPaleGray!72},
  frame=single,
  rulecolor=\color{UTPMuted!35},
  framesep=5pt,
  breaklines=true,
  showstringspaces=false,
  keepspaces=true,
  columns=fullflexible,
  upquote=true,
  aboveskip=0.6em,
  belowskip=0.7em
}

\lstdefinestyle{adscpp}{
  language=C++,
  basicstyle=\ttfamily\scriptsize,
  keywordstyle=\bfseries\color{UTPBlue},
  commentstyle=\itshape\color{UTPMuted},
  stringstyle=\color{UTPGreen!70!black},
  backgroundcolor=\color{UTPPaleGray!72},
  frame=single,
  rulecolor=\color{UTPBlue!28},
  framesep=5pt,
  breaklines=true,
  breakatwhitespace=false,
  showstringspaces=false,
  keepspaces=true,
  columns=fullflexible,
  upquote=true,
  tabsize=4,
  aboveskip=0.7em,
  belowskip=0.8em
}

\lstdefinestyle{mlpython}{
  language=Python,
  basicstyle=\ttfamily\scriptsize,
  keywordstyle=\bfseries\color{UTPBlue},
  commentstyle=\itshape\color{UTPMuted},
  stringstyle=\color{UTPGreen!70!black},
  backgroundcolor=\color{UTPPaleGray!72},
  frame=single,
  rulecolor=\color{UTPBlue!28},
  framesep=5pt,
  breaklines=true,
  breakatwhitespace=false,
  showstringspaces=false,
  keepspaces=true,
  columns=fullflexible,
  upquote=true,
  tabsize=4,
  aboveskip=0.7em,
  belowskip=0.8em
}

\lstdefinestyle{mlmatlab}{
  language=Matlab,
  basicstyle=\ttfamily\scriptsize,
  keywordstyle=\bfseries\color{UTPBlue},
  commentstyle=\itshape\color{UTPMuted},
  stringstyle=\color{UTPGreen!70!black},
  backgroundcolor=\color{UTPPaleGray!72},
  frame=single,
  rulecolor=\color{UTPBlue!28},
  framesep=5pt,
  breaklines=true,
  breakatwhitespace=false,
  showstringspaces=false,
  keepspaces=true,
  columns=fullflexible,
  upquote=true,
  tabsize=4,
  aboveskip=0.7em,
  belowskip=0.8em
}

\lstdefinestyle{mlterminal}{
  basicstyle=\ttfamily\scriptsize,
  backgroundcolor=\color{UTPPaleGray!72},
  frame=single,
  rulecolor=\color{UTPMuted!35},
  framesep=5pt,
  breaklines=true,
  showstringspaces=false,
  keepspaces=true,
  columns=fullflexible,
  upquote=true,
  aboveskip=0.6em,
  belowskip=0.7em
}

  \lstdefinestyle{dcnterminal}{
  basicstyle=\ttfamily\scriptsize,
  backgroundcolor=\color{UTPPaleGray!72},
  frame=single,rulecolor=\color{UTPMuted!35},
  framesep=5pt,
  breaklines=true,
  showstringspaces=false,
  keepspaces=true,
  columns=fullflexible,
  upquote=true,
  aboveskip=0.6em,
  belowskip=0.7em
}

\lstdefinestyle{dsr}{
  language=R,
  basicstyle=\ttfamily\scriptsize,
  keywordstyle=\bfseries\color{UTPBlue},
  commentstyle=\itshape\color{UTPMuted},
  stringstyle=\color{UTPGreen!70!black},
  backgroundcolor=\color{UTPPaleGray!72},
  frame=single,
  rulecolor=\color{UTPBlue!28},
  framesep=5pt,
  breaklines=true,
  breakatwhitespace=false,
  showstringspaces=false,
  keepspaces=true,
  columns=fullflexible,
  upquote=true,
  tabsize=2,
  aboveskip=0.7em,
  belowskip=0.8em
}

\lstdefinestyle{dsterminal}{
  basicstyle=\ttfamily\scriptsize,
  backgroundcolor=\color{UTPPaleGray!72},
  frame=single,
  rulecolor=\color{UTPMuted!35},
  framesep=5pt,
  breaklines=true,
  showstringspaces=false,
  keepspaces=true,
  columns=fullflexible,
  upquote=true,
  aboveskip=0.6em,
  belowskip=0.7em
}

\lstdefinestyle{erpterminal}{
  basicstyle=\ttfamily\scriptsize,
  backgroundcolor=\color{UTPPaleGray!72},
  frame=single,
  rulecolor=\color{UTPBlue!28},
  framesep=5pt,
  breaklines=true,
  showstringspaces=false,
  keepspaces=true,
  columns=fullflexible,
  upquote=true,
  tabsize=4,
  aboveskip=0.7em,
  belowskip=0.8em
}

\lstdefinestyle{mlpython}{
  language=Python,
  basicstyle=\ttfamily\scriptsize,
  keywordstyle=\bfseries\color{UTPBlue},
  commentstyle=\itshape\color{UTPMuted},
  stringstyle=\color{UTPGreen!70!black},
  backgroundcolor=\color{UTPPaleGray!72},
  frame=single,
  rulecolor=\color{UTPBlue!28},
  framesep=5pt,
  breaklines=true,
  breakatwhitespace=false,
  showstringspaces=false,
  keepspaces=true,
  columns=fullflexible,
  upquote=true,
  tabsize=4,
  aboveskip=0.7em,
  belowskip=0.8em
}
\lstdefinestyle{mlterminal}{
  basicstyle=\ttfamily\scriptsize,
  backgroundcolor=\color{UTPPaleGray!72},
  frame=single,
  rulecolor=\color{UTPMuted!35},
  framesep=5pt,
  breaklines=true,
  showstringspaces=false,
  keepspaces=true,
  columns=fullflexible,
  upquote=true,
  aboveskip=0.6em,
  belowskip=0.7em
}

\lstdefinestyle{mlpython}{
  language=Python,
  basicstyle=\ttfamily\scriptsize,
  keywordstyle=\bfseries\color{UTPBlue},
  commentstyle=\itshape\color{UTPMuted},
  stringstyle=\color{UTPGreen!70!black},
  backgroundcolor=\color{UTPPaleGray!72},
  frame=single,
  rulecolor=\color{UTPBlue!28},
  framesep=5pt,
  breaklines=true,
  breakatwhitespace=false,
  showstringspaces=false,
  keepspaces=true,
  columns=fullflexible,
  upquote=true,
  tabsize=4,
  aboveskip=0.7em,
  belowskip=0.8em
}
\lstdefinestyle{mlterminal}{
  basicstyle=\ttfamily\scriptsize,
  backgroundcolor=\color{UTPPaleGray!72},
  frame=single,
  rulecolor=\color{UTPMuted!35},
  framesep=5pt,
  breaklines=true,
  showstringspaces=false,
  keepspaces=true,
  columns=fullflexible,
  upquote=true,
  aboveskip=0.6em,
  belowskip=0.7em
}
\lstdefinestyle{osc}{
  language=C,
  basicstyle=\ttfamily\scriptsize,
  keywordstyle=\bfseries\color{UTPBlue},
  commentstyle=\itshape\color{UTPMuted},
  stringstyle=\color{UTPGreen!70!black},
  backgroundcolor=\color{UTPPaleGray!72},
  frame=single,
  rulecolor=\color{UTPBlue!28},
  framesep=5pt,
  breaklines=true,
  showstringspaces=false,
  keepspaces=true,
  columns=fullflexible,
  upquote=true,
  tabsize=4,
  aboveskip=0.7em,
  belowskip=0.8em
}
\lstdefinestyle{osbash}{
  language=bash,
  basicstyle=\ttfamily\scriptsize,
  keywordstyle=\bfseries\color{UTPBlue},
  commentstyle=\itshape\color{UTPMuted},
  stringstyle=\color{UTPGreen!70!black},
  backgroundcolor=\color{UTPPaleGray!72},
  frame=single,
  rulecolor=\color{UTPMuted!35},
  framesep=5pt,
  breaklines=true,
  showstringspaces=false,
  keepspaces=true,
  columns=fullflexible,
  upquote=true,
  aboveskip=0.6em,
  belowskip=0.7em
}



%      \graphicspath{{../.shared/}}
%
%  Configuration goes BEFORE the \input; overrides go AFTER it.
%  Paths are relative, so compile from inside the course folder.
%
%  Adding a new course? Copy .shared/TEMPLATE-main.tex and follow
%  .shared/ADDING-A-COURSE.md. Every step is a numbered checklist.
%
%  ------------------------------------------------------------------
%  CONTENTS
%  ------------------------------------------------------------------
%    1  Course configuration     the knobs a book may set
%    2  Packages                 the shared package set
%    3  Colour palette           one palette, plus per-course aliases
%    4  Page layout and tables   margins, lists, table styling
%    5  Table of contents        entry formatting for parts/chapters/labs
%    6  Headings                 chapter badge, part pages, sections
%    7  Page furniture           running heads and footers
%    8  Learning boxes           key idea, pitfall, try it, and friends
%    9  Small commands           \term, \cmd, \chapterguide, \mlsteps
%   10  Lab milestone banner     the full-width "Lab N of M" page break
%   11  Stage divider            the part-opening page in longer books
%   12  Listing styles           the base code and terminal styles
% =====================================================================

% --------------------------------------------------------------------
%  1.  COURSE CONFIGURATION
% --------------------------------------------------------------------
%  The three knobs below, plus \utpstyleparts in section 5 and the box
%  titles in section 8, are the whole configuration surface. A book
%  sets what it needs BEFORE the \input and inherits the rest.
%
%    \utprunninghead   the left running head on every page.
%    \utpchapterlabel  the word in the chapter badge. CHAPTER by
%                      default; a book whose chapters are labs sets LAB.
%    \utplabtotal      the N in the "Lab 3 of N" banners.
%
%  \providecommand means "define unless the book already did", which is
%  what makes the \def-before-\input pattern work.

\providecommand{\utprunninghead}{Course Notes}
\providecommand{\utpchapterlabel}{CHAPTER}
\providecommand{\utplabtotal}{10}

% --------------------------------------------------------------------
%  2.  PACKAGES
% --------------------------------------------------------------------
%  The union of what the eight books need. A course that needs
%  something extra loads it in its own file, under the heading
%  "packages only this book needs" (algorithm2e, caption, float and
%  multirow are the only current examples).
%
%  The tikz library list is the union across all books, so a diagram
%  copied from one course compiles in another.

\usepackage[T1]{fontenc}
\IfFileExists{glyphtounicode.tex}{\input{glyphtounicode}\pdfgentounicode=1}{}
\usepackage[american]{babel}
\IfFileExists{sourcesanspro.sty}
  {\usepackage[default]{sourcesanspro}}
  {\usepackage[scaled]{helvet}\renewcommand{\familydefault}{\sfdefault}}
\DeclareFontShape{T1}{cmss}{b}{n}{<->ssub*cmss/bx/n}{}
\DeclareFontShape{TS1}{cmss}{b}{n}{<->ssub*cmss/bx/n}{}
\usepackage[margin=21mm,headheight=22pt]{geometry}
\usepackage{amsmath,amssymb,mathtools,bm}
\usepackage{microtype}
\usepackage{enumitem}
\usepackage{booktabs}
\usepackage{tabularx}
\usepackage{array}
\usepackage{ragged2e}
\usepackage{etoolbox}
\usepackage{titlesec}
\usepackage{tocloft}
\usepackage{fancyhdr}
\usepackage[table]{xcolor}
\usepackage{tcolorbox}
\tcbuselibrary{breakable,skins}
\usepackage[hidelinks]{hyperref}
\usepackage{bookmark}
\usepackage{listings}
\usepackage{graphicx}
\usepackage{tikz}
\usetikzlibrary{angles,arrows.meta,backgrounds,calc,chains,decorations.pathmorphing,decorations.pathreplacing,fit,matrix,patterns,positioning,quotes,shadows.blur,shapes.geometric,shapes.misc}
\usepackage{pgfplots}

% --------------------------------------------------------------------
%  3.  COLOUR PALETTE
% --------------------------------------------------------------------
%  One palette, defined once under canonical UTP* names. Use these
%  in any new material.
%
%  The existing chapters were written against course-prefixed names
%  (MLNavy, ADSPaleBlue, ...) in over two thousand places, so rather
%  than rewrite them the block below aliases each prefix onto the
%  canonical colours. All four sets are always defined, so a chapter
%  moved between courses still resolves.

\definecolor{UTPNavy}{HTML}{18324A}
\definecolor{UTPBlue}{HTML}{246B9E}
\definecolor{UTPTeal}{HTML}{138A80}
\definecolor{UTPGreen}{HTML}{2D7D46}
\definecolor{UTPOrange}{HTML}{C86B24}
\definecolor{UTPRed}{HTML}{B34843}
\definecolor{UTPGold}{HTML}{D4A13E}
\definecolor{UTPInk}{HTML}{263238}
\definecolor{UTPMuted}{HTML}{60717C}
\definecolor{UTPPaleBlue}{HTML}{EEF6FB}
\definecolor{UTPPaleTeal}{HTML}{EEF8F6}
\definecolor{UTPPaleOrange}{HTML}{FFF5EA}
\definecolor{UTPPaleRed}{HTML}{FFF0EF}
\definecolor{UTPPaleGray}{HTML}{F4F6F7}

% Course files were written against prefixed colour names. Rather than edit
% 2000+ references across the chapters, alias each prefix onto the canonical
% palette. Only these four prefixes appear anywhere in the books.
\def\utpmkalias#1{%
  \colorlet{#1Navy}{UTPNavy}\colorlet{#1Blue}{UTPBlue}%
  \colorlet{#1Teal}{UTPTeal}\colorlet{#1Green}{UTPGreen}%
  \colorlet{#1Orange}{UTPOrange}\colorlet{#1Red}{UTPRed}%
  \colorlet{#1Ink}{UTPInk}\colorlet{#1Muted}{UTPMuted}%
  \colorlet{#1PaleBlue}{UTPPaleBlue}\colorlet{#1PaleTeal}{UTPPaleTeal}%
  \colorlet{#1PaleOrange}{UTPPaleOrange}\colorlet{#1PaleRed}{UTPPaleRed}%
  \colorlet{#1PaleGray}{UTPPaleGray}%
}
\utpmkalias{ML}
\utpmkalias{ADS}
\utpmkalias{DS}
\utpmkalias{ERP}
\color{UTPInk}

% --------------------------------------------------------------------
%  4.  PAGE LAYOUT AND TABLES
% --------------------------------------------------------------------
%  Paragraph spacing, list indents, and the table look: alternating
%  row tint, top-aligned paragraph columns, ragged-right cells.
%
%  A book typeset tighter overrides \arraystretch, \tabcolsep or
%  \emergencystretch after the input. Data Communication is the one
%  that currently does.

\setlength{\parindent}{0pt}
\setlength{\parskip}{0.55em}
\setlist[itemize]{leftmargin=1.45em,itemsep=0.22em,topsep=0.25em}
\setlist[enumerate]{leftmargin=1.75em,itemsep=0.28em,topsep=0.3em}
\setlist[description]{font=\bfseries\color{UTPNavy},itemsep=0.25em,topsep=0.3em}
\renewcommand{\arraystretch}{1.34}
\setlength{\tabcolsep}{6.5pt}
\arrayrulecolor{UTPBlue!45}
% Keep paragraph-style columns top-aligned and ragged. Tables mix fixed p
% columns with flexible X columns; giving X the same baseline stops one
% column floating halfway down a multi-line row.
\renewcommand{\tabularxcolumn}[1]{p{#1}}
\makeatletter
\g@addto@macro\@arrayparboxrestore{\RaggedRight\arraybackslash}
\makeatother
\AtBeginEnvironment{tabular}{\small\setlength{\extrarowheight}{0.6pt}\rowcolors{2}{UTPPaleBlue!72}{white}}
\AtBeginEnvironment{tabularx}{\small\setlength{\extrarowheight}{0.6pt}\rowcolors{2}{UTPPaleBlue!72}{white}}
\AtBeginEnvironment{figure}{\centering}
\emergencystretch=1.5em
\hbadness=10000
\hfuzz=8pt

\pgfplotsset{
  compat=1.18,
  every axis/.append style={
    axis line style={draw=UTPMuted!85,line width=0.55pt},
    tick style={draw=UTPMuted!85},
    tick label style={font=\scriptsize\color{UTPMuted}},
    label style={font=\small\color{UTPNavy}},
    title style={font=\small\bfseries\color{UTPNavy}},
    legend style={font=\scriptsize,draw=UTPBlue!22,fill=white,rounded corners=1pt}
  },
  every axis plot/.append style={line width=1.15pt}
}
\tikzset{every picture/.append style={line cap=round,line join=round}}

% --------------------------------------------------------------------
%  5.  TABLE OF CONTENTS
% --------------------------------------------------------------------
%  Formatting for part, chapter and lab entries. The lab level is a
%  custom tocloft list so lab banners appear in the contents.
%
%  Part styling is skipped when \utpstyleparts is 0, which is what a
%  book without \part divisions should set.

\addto\captionsamerican{\renewcommand{\contentsname}{Learning Path}}
\newlistentry{lab}{toc}{-1}
\providecommand{\utpstyleparts}{1}
\ifnum\utpstyleparts=1
\renewcommand{\cftpartpresnum}{}
\renewcommand{\cftpartaftersnum}{}
\setlength{\cftpartnumwidth}{5.2em}
\setlength{\cftbeforepartskip}{0.38em}
\renewcommand{\cftpartfont}{\normalsize\bfseries\color{UTPNavy}}
\renewcommand{\cftpartpagefont}{\normalsize\bfseries\color{UTPNavy}}
\renewcommand{\cftpartleader}{\hfill}
\fi
\renewcommand{\cftchappresnum}{Chapter~}
\renewcommand{\cftchapaftersnum}{\quad}
\setlength{\cftchapindent}{1.5em}
\setlength{\cftchapnumwidth}{6em}
\setlength{\cftbeforechapskip}{0pt}
\renewcommand{\cftchapfont}{\normalfont}
\renewcommand{\cftchappagefont}{\normalfont}
\renewcommand{\cftchapleader}{\cftdotfill{\cftdotsep}}
\setlength{\cftlabindent}{3em}
\setlength{\cftlabnumwidth}{0em}
\setlength{\cftbeforelabskip}{0.10em}
\renewcommand{\cftlabfont}{\bfseries\color{UTPOrange!90!black}}
\renewcommand{\cftlabpagefont}{\bfseries\color{UTPOrange!90!black}}
\renewcommand{\cftlableader}{\color{UTPOrange!65!UTPMuted}\cftdotfill{\cftlabdotsep}}

% --------------------------------------------------------------------
%  6.  HEADINGS
% --------------------------------------------------------------------
%  The chapter badge (a coloured block with a word and the number),
%  part pages, and section spacing.
%
%  The word in the badge is \utpchapterlabel; Data Communication sets
%  it to LAB because each of its chapters is a lab.

\titleformat{\chapter}[display]
  {\bfseries\color{UTPNavy}}
  {\filleft\colorbox{UTPBlue}{\parbox{2.55cm}{\centering\color{white}\Large
    \utpchapterlabel\\[-0.15em]\Huge\thechapter}}}
  {0.65em}{\Huge\raggedright}
  [\vspace{0.25em}{\color{UTPBlue}\titlerule[1.2pt]}]
\renewcommand{\thepart}{Stage~\arabic{part}}
\titleformat{\part}[display]
  {\centering\color{UTPNavy}\bfseries}
  {\Large\color{UTPTeal}\thepart}{0.45em}{\Huge}
  [\vspace{0.55em}\color{UTPTeal}\rule{0.45\linewidth}{1.5pt}]
\titleformat{\section}{\Large\bfseries\color{UTPNavy}}{\thesection}{0.55em}{}
\titleformat{\subsection}{\large\bfseries\color{UTPBlue}}{\thesubsection}{0.5em}{}
\titlespacing*{\chapter}{0pt}{-1.8em}{1.0em}
\titlespacing*{\part}{0pt}{0pt}{1em}
\titlespacing*{\section}{0pt}{1.2em}{0.35em}
\titlespacing*{\subsection}{0pt}{0.95em}{0.2em}

% --------------------------------------------------------------------
%  7.  PAGE FURNITURE
% --------------------------------------------------------------------
%  Running heads and footers. The left header is \utprunninghead;
%  the right one follows the current chapter.

\pagestyle{fancy}
\fancyhf{}
\lhead{\small\color{UTPMuted}\utprunninghead}
\rhead{\small\color{UTPMuted}\nouppercase{\leftmark}}
\cfoot{\small\color{UTPMuted}\thepage}
\renewcommand{\headrulewidth}{0.5pt}
\renewcommand{\headrule}{\hbox to\headwidth{\color{UTPPaleBlue}\leaders\hrule height \headrulewidth\hfill}}
\renewcommand{\chaptermark}[1]{\markboth{#1}{}}
\fancypagestyle{plain}{%
  \fancyhf{}%
  \cfoot{\small\color{UTPMuted}\thepage}%
  \renewcommand{\headrulewidth}{0pt}%
}

% --------------------------------------------------------------------
%  8.  LEARNING BOXES
% --------------------------------------------------------------------
%  The coloured callouts the chapters use. Every title is a macro, so
%  a book can rename a box without restating its geometry:
%
%      \def\utptitlepitfall{Technical note}   % before the input
%
%  To change how a box LOOKS in one book, use \renewtcolorbox after
%  the input instead. To add a box only one book needs, declare it
%  there with \newtcolorbox.

% Box titles -- override any of these before the input line.
\providecommand{\utptitlekeyidea}{Key idea}
\providecommand{\utptitleworked}{Worked example}
\providecommand{\utptitlepitfall}{Common mistake}
\providecommand{\utptitletryit}{Try it yourself}
\providecommand{\utptitlequickcheck}{Check your understanding}
\providecommand{\utptitlerecap}{Chapter recap}
\providecommand{\utptitledeliverable}{Lab deliverable}
\providecommand{\utptitlelabmode}{How much is scaffolded?}
\providecommand{\utptitleintuition}{Plain idea}
\providecommand{\utptitlefirstread}{Core path}
\providecommand{\utptitleglance}{Chapter at a glance}
\providecommand{\utptitledebug}{Debugging task}

% Explaining an idea
\newtcolorbox{keyidea}{enhanced,breakable,colback=UTPPaleTeal,colframe=UTPTeal, boxrule=0.8pt,arc=2mm,left=2.5mm,right=2.5mm,top=1.8mm,bottom=1.8mm, fonttitle=\bfseries\color{UTPNavy},title=\utptitlekeyidea}
\newtcolorbox{workedexample}{enhanced,breakable,colback=UTPPaleBlue,colframe=UTPBlue, boxrule=0.8pt,arc=2mm,left=2.5mm,right=2.5mm,top=1.8mm,bottom=1.8mm, fonttitle=\bfseries\color{UTPNavy},title=\utptitleworked}
\newtcolorbox{pitfall}{enhanced,breakable,colback=UTPPaleRed,colframe=UTPRed, boxrule=0.8pt,arc=2mm,left=2.5mm,right=2.5mm,top=1.8mm,bottom=1.8mm, fonttitle=\bfseries\color{UTPNavy},title=\utptitlepitfall}
% Asking the reader to do something
\newtcolorbox{tryit}{enhanced,breakable,colback=white,colframe=UTPTeal, boxrule=0.9pt,arc=2mm,left=2.5mm,right=2.5mm,top=1.8mm,bottom=1.8mm, fonttitle=\bfseries\color{white},colbacktitle=UTPTeal,title=\utptitletryit}
\newtcolorbox{quickcheck}{enhanced,breakable,colback=white,colframe=UTPOrange, boxrule=0.9pt,arc=2mm,left=2.5mm,right=2.5mm,top=1.8mm,bottom=1.8mm, fonttitle=\bfseries\color{white},colbacktitle=UTPOrange,title=\utptitlequickcheck}
\newtcolorbox{chapterrecap}{enhanced,breakable,colback=UTPPaleTeal,colframe=UTPGreen, boxrule=0.9pt,arc=2mm,left=2.5mm,right=2.5mm,top=1.8mm,bottom=1.8mm, fonttitle=\bfseries\color{white},colbacktitle=UTPGreen,title=\utptitlerecap}
% Lab structure
\newtcolorbox{labdeliverable}{enhanced,breakable,colback=white,colframe=UTPGreen, boxrule=0.9pt,arc=2mm,left=2.6mm,right=2.6mm,top=1.9mm,bottom=1.9mm, fonttitle=\bfseries\color{white},colbacktitle=UTPGreen,title=\utptitledeliverable}
\newtcolorbox{labmode}{enhanced,breakable,colback=UTPPaleBlue!52,colframe=UTPBlue, boxrule=0.75pt,arc=2mm,left=2.6mm,right=2.6mm,top=1.7mm,bottom=1.7mm, fonttitle=\bfseries\color{white},colbacktitle=UTPBlue,title=\utptitlelabmode}
% Softer, lighter-weight asides
\newtcolorbox{intuition}{enhanced,breakable,colback=white,colframe=UTPOrange!42, boxrule=0.5pt,arc=2mm,left=2.5mm,right=2.5mm,top=1.5mm,bottom=1.5mm, fonttitle=\bfseries\color{UTPOrange!85!black},title=\utptitleintuition}
\newtcolorbox{firstread}{enhanced,breakable,colback=white,colframe=UTPBlue!42, boxrule=0.5pt,arc=2mm,left=2.5mm,right=2.5mm,top=1.5mm,bottom=1.5mm, fonttitle=\bfseries\color{UTPBlue}, title=\utptitlefirstread}
\newtcolorbox{debugtask}{enhanced,breakable,colback=UTPPaleRed!58,colframe=UTPRed, boxrule=0.85pt,arc=2mm,left=2.6mm,right=2.6mm,top=1.9mm,bottom=1.9mm, fonttitle=\bfseries\color{white},colbacktitle=UTPRed,title=\utptitledebug}
\newtcolorbox{analogy}{enhanced,breakable,colback=white,colframe=UTPOrange!38, boxrule=0.5pt,arc=2mm,left=2.5mm,right=2.5mm,top=1.5mm,bottom=1.5mm, fonttitle=\bfseries\color{UTPOrange!85!black},title=Analogy}
\newtcolorbox{coursechoice}{enhanced,breakable,colback=white,colframe=UTPNavy, boxrule=0.8pt,arc=2mm,left=2.8mm,right=2.8mm,top=2mm,bottom=2mm, fonttitle=\bfseries\color{white},colbacktitle=UTPNavy,title=Course choice}
\newtcolorbox{decisionmemo}{enhanced,breakable,colback=UTPPaleOrange!55,colframe=UTPOrange, boxrule=0.85pt,arc=2mm,left=2.6mm,right=2.6mm,top=1.9mm,bottom=1.9mm, fonttitle=\bfseries\color{white},colbacktitle=UTPOrange,title=Decision memo}
\newtcolorbox{libnote}{enhanced,breakable,colback=white,colframe=UTPMuted!48, boxrule=0.5pt,arc=2mm,left=2.5mm,right=2.5mm,top=1.5mm,bottom=1.5mm, fonttitle=\bfseries\color{UTPNavy},title=Reference note}
\newtcolorbox{miniproject}{enhanced,breakable,colback=UTPPaleTeal!52,colframe=UTPGreen, boxrule=0.9pt,arc=2mm,left=2.6mm,right=2.6mm,top=1.9mm,bottom=1.9mm, fonttitle=\bfseries\color{white},colbacktitle=UTPGreen,title=Mini-project}
% Boxes that take a title argument
\newtcolorbox{modelcard}[1]{enhanced,breakable,colback=white,colframe=UTPBlue, boxrule=0.9pt,arc=2mm,left=2.5mm,right=2.5mm,top=1.8mm,bottom=1.8mm, fonttitle=\bfseries\color{white},colbacktitle=UTPBlue,title=#1}
\newtcolorbox{optionaltopic}{enhanced,breakable,colback=white,colframe=UTPMuted!58, boxrule=0.5pt,arc=2mm,left=2.5mm,right=2.5mm,top=1.5mm,bottom=1.5mm, fonttitle=\bfseries\color{UTPMuted},title=Beyond the core}
\newtcolorbox{secondpass}{enhanced,breakable,colback=UTPPaleGray!82,colframe=UTPMuted, boxrule=0.75pt,arc=2mm,left=2.6mm,right=2.6mm,top=1.9mm,bottom=1.9mm, fonttitle=\bfseries\color{white},colbacktitle=UTPMuted,title=Second pass}
\newtcolorbox{labproject}[1]{enhanced,breakable,colback=white,colframe=UTPBlue, boxrule=0.9pt,arc=2mm,left=2.6mm,right=2.6mm,top=1.9mm,bottom=1.9mm, fonttitle=\bfseries\color{white},colbacktitle=UTPBlue,title={#1}}
% Frames for a figure or a table
\newtcolorbox{mlfigure}{enhanced,width=\linewidth,colback=white,colframe=UTPBlue!28, boxrule=0.7pt,arc=2mm, left=2.2mm,right=2.2mm,top=2.2mm,bottom=1.8mm, before skip=0.75em,after skip=0.85em, before upper={\centering}, halign=center}
\newtcolorbox{mltable}{enhanced,width=\linewidth,colback=white,colframe=UTPTeal!32, boxrule=0.7pt,arc=2mm, left=2.2mm,right=2.2mm,top=1.8mm,bottom=1.8mm, before skip=0.75em,after skip=0.85em, halign=center}

% --------------------------------------------------------------------
%  9.  SMALL COMMANDS
% --------------------------------------------------------------------
%  \term for a defined term, \cmd for inline code, \chapref and
%  \secref for cross-references, and the mlsteps numbered-step list.
%
%      \term{overfitting}          a term being defined
%      \cmd{git status}            inline code
%      \chapref{ch:trees}          "Chapter 12"
%      \secref{sec:pruning}        "Section 12.3"
%      \chapterguide{prior}{outcomes}   the box that opens a chapter
%      \begin{mlsteps} ... \end{mlsteps}  a numbered Step 1/Step 2 list

\newcommand{\term}[1]{\textbf{\color{UTPNavy}#1}}
\newcommand{\cmd}[1]{\texttt{#1}}
\newcommand{\chapref}[1]{Chapter~\ref{#1}}
\newcommand{\secref}[1]{Section~\ref{#1}}
\newcommand{\chapterguide}[2]{%
  \begin{tcolorbox}[enhanced,breakable,colback=white,colframe=UTPNavy,boxrule=0.8pt,
    arc=2mm,left=2.5mm,right=2.5mm,top=2mm,bottom=2mm,
    title=\utptitleglance,fonttitle=\bfseries\color{white},colbacktitle=UTPNavy]
    \textbf{Before you start:} #1\par
    \textbf{By the end, you can:} #2
  \end{tcolorbox}}
\newenvironment{mlsteps}
  {\begin{enumerate}[label=\textcolor{UTPBlue}{\bfseries Step~\arabic*:},leftmargin=4.5em,itemsep=0.38em]}
  {\end{enumerate}}

% --------------------------------------------------------------------
%  10.  LAB MILESTONE BANNER
% --------------------------------------------------------------------
%  The full-width page break that opens each lab. Called from the lab
%  files as \labsection{number}{title}{label}.
%
%  The "of N" comes from \utplabtotal, so a book states its lab count
%  once in its configuration rather than in every lab file.

\newcommand{\labsection}[3]{%
  \clearpage
  \phantomsection
  \addcontentsline{toc}{lab}{Lab #1 of \utplabtotal\space--- #2}%
  \markboth{Lab #1 of \utplabtotal\space--- #2}{}%
  \begin{tcolorbox}[enhanced,width=\linewidth,colback=UTPNavy,colframe=UTPNavy,
    boxrule=0pt,arc=2mm,left=5mm,right=5mm,top=5mm,bottom=5mm,
    before skip=0pt,after skip=1.2em]
    {\large\bfseries\color{UTPTeal!35!white} PRACTICAL MILESTONE\par}
    \vspace{0.25em}
    {\Huge\bfseries\color{white} Lab #1 of \utplabtotal\par}
    \vspace{0.20em}
    {\Large\bfseries\color{white} #2\par}
  \end{tcolorbox}
  \label{#3}%
}

% --------------------------------------------------------------------
%  11.  STAGE DIVIDER
% --------------------------------------------------------------------
%  A full-page divider used by the longer books to open a part:
%  \substage{3A}{Supervised Workflow}{one-line summary}.
%  Harmless to leave unused.

\newcommand{\substage}[3]{%
  \cleardoublepage
  \thispagestyle{empty}
  \phantomsection
  \addcontentsline{toc}{chapter}{Stage #1 --- #2}
  \begin{center}
    \vspace*{0.20\textheight}
    {\Large\bfseries\color{UTPTeal} Stage #1\par}
    \vspace{0.45em}
    {\Huge\bfseries\color{UTPNavy} #2\par}
    \vspace{0.65em}
    {\color{UTPTeal}\rule{0.46\linewidth}{1.5pt}\par}
    \vspace{1.1em}
    \begin{minipage}{0.76\linewidth}
      \centering\large\color{UTPMuted} #3
    \end{minipage}
  \end{center}
  \cleardoublepage
}

% --------------------------------------------------------------------
%  12.  LISTING STYLES
% --------------------------------------------------------------------
%  Two base styles. A course builds its own on top of them so the
%  colours and frame stay consistent:
%
%      \lstdefinestyle{mlpython}{style=utpcode,language=Python}
%      \lstdefinestyle{mlterminal}{style=utpterminal}

\lstdefinestyle{utpcode}{
  basicstyle=\ttfamily\scriptsize,
  keywordstyle=\bfseries\color{UTPBlue},
  commentstyle=\itshape\color{UTPMuted},
  stringstyle=\color{UTPGreen!70!black},
  backgroundcolor=\color{UTPPaleGray!72},
  frame=single,
  rulecolor=\color{UTPBlue!28},
  framesep=5pt,
  breaklines=true,
  breakatwhitespace=false,
  showstringspaces=false,
  keepspaces=true,
  columns=fullflexible,
  upquote=true,
  tabsize=4,
  aboveskip=0.7em,
  belowskip=0.8em
}
\lstdefinestyle{utpterminal}{
  basicstyle=\ttfamily\scriptsize,
  backgroundcolor=\color{UTPPaleGray!72},
  frame=single,
  rulecolor=\color{UTPMuted!35},
  framesep=5pt,
  breaklines=true,
  showstringspaces=false,
  keepspaces=true,
  columns=fullflexible,
  upquote=true,
  aboveskip=0.6em,
  belowskip=0.7em
}

\lstdefinestyle{adscpp}{
  language=C++,
  basicstyle=\ttfamily\scriptsize,
  keywordstyle=\bfseries\color{UTPBlue},
  commentstyle=\itshape\color{UTPMuted},
  stringstyle=\color{UTPGreen!70!black},
  backgroundcolor=\color{UTPPaleGray!72},
  frame=single,
  rulecolor=\color{UTPBlue!28},
  framesep=5pt,
  breaklines=true,
  breakatwhitespace=false,
  showstringspaces=false,
  keepspaces=true,
  columns=fullflexible,
  upquote=true,
  tabsize=4,
  aboveskip=0.7em,
  belowskip=0.8em
}

\lstdefinestyle{mlpython}{
  language=Python,
  basicstyle=\ttfamily\scriptsize,
  keywordstyle=\bfseries\color{UTPBlue},
  commentstyle=\itshape\color{UTPMuted},
  stringstyle=\color{UTPGreen!70!black},
  backgroundcolor=\color{UTPPaleGray!72},
  frame=single,
  rulecolor=\color{UTPBlue!28},
  framesep=5pt,
  breaklines=true,
  breakatwhitespace=false,
  showstringspaces=false,
  keepspaces=true,
  columns=fullflexible,
  upquote=true,
  tabsize=4,
  aboveskip=0.7em,
  belowskip=0.8em
}

\lstdefinestyle{mlmatlab}{
  language=Matlab,
  basicstyle=\ttfamily\scriptsize,
  keywordstyle=\bfseries\color{UTPBlue},
  commentstyle=\itshape\color{UTPMuted},
  stringstyle=\color{UTPGreen!70!black},
  backgroundcolor=\color{UTPPaleGray!72},
  frame=single,
  rulecolor=\color{UTPBlue!28},
  framesep=5pt,
  breaklines=true,
  breakatwhitespace=false,
  showstringspaces=false,
  keepspaces=true,
  columns=fullflexible,
  upquote=true,
  tabsize=4,
  aboveskip=0.7em,
  belowskip=0.8em
}

\lstdefinestyle{mlterminal}{
  basicstyle=\ttfamily\scriptsize,
  backgroundcolor=\color{UTPPaleGray!72},
  frame=single,
  rulecolor=\color{UTPMuted!35},
  framesep=5pt,
  breaklines=true,
  showstringspaces=false,
  keepspaces=true,
  columns=fullflexible,
  upquote=true,
  aboveskip=0.6em,
  belowskip=0.7em
}

  \lstdefinestyle{dcnterminal}{
  basicstyle=\ttfamily\scriptsize,
  backgroundcolor=\color{UTPPaleGray!72},
  frame=single,rulecolor=\color{UTPMuted!35},
  framesep=5pt,
  breaklines=true,
  showstringspaces=false,
  keepspaces=true,
  columns=fullflexible,
  upquote=true,
  aboveskip=0.6em,
  belowskip=0.7em
}

\lstdefinestyle{dsr}{
  language=R,
  basicstyle=\ttfamily\scriptsize,
  keywordstyle=\bfseries\color{UTPBlue},
  commentstyle=\itshape\color{UTPMuted},
  stringstyle=\color{UTPGreen!70!black},
  backgroundcolor=\color{UTPPaleGray!72},
  frame=single,
  rulecolor=\color{UTPBlue!28},
  framesep=5pt,
  breaklines=true,
  breakatwhitespace=false,
  showstringspaces=false,
  keepspaces=true,
  columns=fullflexible,
  upquote=true,
  tabsize=2,
  aboveskip=0.7em,
  belowskip=0.8em
}

\lstdefinestyle{dsterminal}{
  basicstyle=\ttfamily\scriptsize,
  backgroundcolor=\color{UTPPaleGray!72},
  frame=single,
  rulecolor=\color{UTPMuted!35},
  framesep=5pt,
  breaklines=true,
  showstringspaces=false,
  keepspaces=true,
  columns=fullflexible,
  upquote=true,
  aboveskip=0.6em,
  belowskip=0.7em
}

\lstdefinestyle{erpterminal}{
  basicstyle=\ttfamily\scriptsize,
  backgroundcolor=\color{UTPPaleGray!72},
  frame=single,
  rulecolor=\color{UTPBlue!28},
  framesep=5pt,
  breaklines=true,
  showstringspaces=false,
  keepspaces=true,
  columns=fullflexible,
  upquote=true,
  tabsize=4,
  aboveskip=0.7em,
  belowskip=0.8em
}

\lstdefinestyle{mlpython}{
  language=Python,
  basicstyle=\ttfamily\scriptsize,
  keywordstyle=\bfseries\color{UTPBlue},
  commentstyle=\itshape\color{UTPMuted},
  stringstyle=\color{UTPGreen!70!black},
  backgroundcolor=\color{UTPPaleGray!72},
  frame=single,
  rulecolor=\color{UTPBlue!28},
  framesep=5pt,
  breaklines=true,
  breakatwhitespace=false,
  showstringspaces=false,
  keepspaces=true,
  columns=fullflexible,
  upquote=true,
  tabsize=4,
  aboveskip=0.7em,
  belowskip=0.8em
}
\lstdefinestyle{mlterminal}{
  basicstyle=\ttfamily\scriptsize,
  backgroundcolor=\color{UTPPaleGray!72},
  frame=single,
  rulecolor=\color{UTPMuted!35},
  framesep=5pt,
  breaklines=true,
  showstringspaces=false,
  keepspaces=true,
  columns=fullflexible,
  upquote=true,
  aboveskip=0.6em,
  belowskip=0.7em
}

\lstdefinestyle{mlpython}{
  language=Python,
  basicstyle=\ttfamily\scriptsize,
  keywordstyle=\bfseries\color{UTPBlue},
  commentstyle=\itshape\color{UTPMuted},
  stringstyle=\color{UTPGreen!70!black},
  backgroundcolor=\color{UTPPaleGray!72},
  frame=single,
  rulecolor=\color{UTPBlue!28},
  framesep=5pt,
  breaklines=true,
  breakatwhitespace=false,
  showstringspaces=false,
  keepspaces=true,
  columns=fullflexible,
  upquote=true,
  tabsize=4,
  aboveskip=0.7em,
  belowskip=0.8em
}
\lstdefinestyle{mlterminal}{
  basicstyle=\ttfamily\scriptsize,
  backgroundcolor=\color{UTPPaleGray!72},
  frame=single,
  rulecolor=\color{UTPMuted!35},
  framesep=5pt,
  breaklines=true,
  showstringspaces=false,
  keepspaces=true,
  columns=fullflexible,
  upquote=true,
  aboveskip=0.6em,
  belowskip=0.7em
}
\lstdefinestyle{osc}{
  language=C,
  basicstyle=\ttfamily\scriptsize,
  keywordstyle=\bfseries\color{UTPBlue},
  commentstyle=\itshape\color{UTPMuted},
  stringstyle=\color{UTPGreen!70!black},
  backgroundcolor=\color{UTPPaleGray!72},
  frame=single,
  rulecolor=\color{UTPBlue!28},
  framesep=5pt,
  breaklines=true,
  showstringspaces=false,
  keepspaces=true,
  columns=fullflexible,
  upquote=true,
  tabsize=4,
  aboveskip=0.7em,
  belowskip=0.8em
}
\lstdefinestyle{osbash}{
  language=bash,
  basicstyle=\ttfamily\scriptsize,
  keywordstyle=\bfseries\color{UTPBlue},
  commentstyle=\itshape\color{UTPMuted},
  stringstyle=\color{UTPGreen!70!black},
  backgroundcolor=\color{UTPPaleGray!72},
  frame=single,
  rulecolor=\color{UTPMuted!35},
  framesep=5pt,
  breaklines=true,
  showstringspaces=false,
  keepspaces=true,
  columns=fullflexible,
  upquote=true,
  aboveskip=0.6em,
  belowskip=0.7em
}



%      \graphicspath{{../.shared/}}
%
%  Configuration goes BEFORE the \input; overrides go AFTER it.
%  Paths are relative, so compile from inside the course folder.
%
%  Adding a new course? Copy .shared/TEMPLATE-main.tex and follow
%  .shared/ADDING-A-COURSE.md. Every step is a numbered checklist.
%
%  ------------------------------------------------------------------
%  CONTENTS
%  ------------------------------------------------------------------
%    1  Course configuration     the knobs a book may set
%    2  Packages                 the shared package set
%    3  Colour palette           one palette, plus per-course aliases
%    4  Page layout and tables   margins, lists, table styling
%    5  Table of contents        entry formatting for parts/chapters/labs
%    6  Headings                 chapter badge, part pages, sections
%    7  Page furniture           running heads and footers
%    8  Learning boxes           key idea, pitfall, try it, and friends
%    9  Small commands           \term, \cmd, \chapterguide, \mlsteps
%   10  Lab milestone banner     the full-width "Lab N of M" page break
%   11  Stage divider            the part-opening page in longer books
%   12  Listing styles           the base code and terminal styles
% =====================================================================

% --------------------------------------------------------------------
%  1.  COURSE CONFIGURATION
% --------------------------------------------------------------------
%  The three knobs below, plus \utpstyleparts in section 5 and the box
%  titles in section 8, are the whole configuration surface. A book
%  sets what it needs BEFORE the \input and inherits the rest.
%
%    \utprunninghead   the left running head on every page.
%    \utpchapterlabel  the word in the chapter badge. CHAPTER by
%                      default; a book whose chapters are labs sets LAB.
%    \utplabtotal      the N in the "Lab 3 of N" banners.
%
%  \providecommand means "define unless the book already did", which is
%  what makes the \def-before-\input pattern work.

\providecommand{\utprunninghead}{Course Notes}
\providecommand{\utpchapterlabel}{CHAPTER}
\providecommand{\utplabtotal}{10}

% --------------------------------------------------------------------
%  2.  PACKAGES
% --------------------------------------------------------------------
%  The union of what the eight books need. A course that needs
%  something extra loads it in its own file, under the heading
%  "packages only this book needs" (algorithm2e, caption, float and
%  multirow are the only current examples).
%
%  The tikz library list is the union across all books, so a diagram
%  copied from one course compiles in another.

\usepackage[T1]{fontenc}
\IfFileExists{glyphtounicode.tex}{\input{glyphtounicode}\pdfgentounicode=1}{}
\usepackage[american]{babel}
\IfFileExists{sourcesanspro.sty}
  {\usepackage[default]{sourcesanspro}}
  {\usepackage[scaled]{helvet}\renewcommand{\familydefault}{\sfdefault}}
\DeclareFontShape{T1}{cmss}{b}{n}{<->ssub*cmss/bx/n}{}
\DeclareFontShape{TS1}{cmss}{b}{n}{<->ssub*cmss/bx/n}{}
\usepackage[margin=21mm,headheight=22pt]{geometry}
\usepackage{amsmath,amssymb,mathtools,bm}
\usepackage{microtype}
\usepackage{enumitem}
\usepackage{booktabs}
\usepackage{tabularx}
\usepackage{array}
\usepackage{ragged2e}
\usepackage{etoolbox}
\usepackage{titlesec}
\usepackage{tocloft}
\usepackage{fancyhdr}
\usepackage[table]{xcolor}
\usepackage{tcolorbox}
\tcbuselibrary{breakable,skins}
\usepackage[hidelinks]{hyperref}
\usepackage{bookmark}
\usepackage{listings}
\usepackage{graphicx}
\usepackage{tikz}
\usetikzlibrary{angles,arrows.meta,backgrounds,calc,chains,decorations.pathmorphing,decorations.pathreplacing,fit,matrix,patterns,positioning,quotes,shadows.blur,shapes.geometric,shapes.misc}
\usepackage{pgfplots}

% --------------------------------------------------------------------
%  3.  COLOUR PALETTE
% --------------------------------------------------------------------
%  One palette, defined once under canonical UTP* names. Use these
%  in any new material.
%
%  The existing chapters were written against course-prefixed names
%  (MLNavy, ADSPaleBlue, ...) in over two thousand places, so rather
%  than rewrite them the block below aliases each prefix onto the
%  canonical colours. All four sets are always defined, so a chapter
%  moved between courses still resolves.

\definecolor{UTPNavy}{HTML}{18324A}
\definecolor{UTPBlue}{HTML}{246B9E}
\definecolor{UTPTeal}{HTML}{138A80}
\definecolor{UTPGreen}{HTML}{2D7D46}
\definecolor{UTPOrange}{HTML}{C86B24}
\definecolor{UTPRed}{HTML}{B34843}
\definecolor{UTPGold}{HTML}{D4A13E}
\definecolor{UTPInk}{HTML}{263238}
\definecolor{UTPMuted}{HTML}{60717C}
\definecolor{UTPPaleBlue}{HTML}{EEF6FB}
\definecolor{UTPPaleTeal}{HTML}{EEF8F6}
\definecolor{UTPPaleOrange}{HTML}{FFF5EA}
\definecolor{UTPPaleRed}{HTML}{FFF0EF}
\definecolor{UTPPaleGray}{HTML}{F4F6F7}

% Course files were written against prefixed colour names. Rather than edit
% 2000+ references across the chapters, alias each prefix onto the canonical
% palette. Only these four prefixes appear anywhere in the books.
\def\utpmkalias#1{%
  \colorlet{#1Navy}{UTPNavy}\colorlet{#1Blue}{UTPBlue}%
  \colorlet{#1Teal}{UTPTeal}\colorlet{#1Green}{UTPGreen}%
  \colorlet{#1Orange}{UTPOrange}\colorlet{#1Red}{UTPRed}%
  \colorlet{#1Ink}{UTPInk}\colorlet{#1Muted}{UTPMuted}%
  \colorlet{#1PaleBlue}{UTPPaleBlue}\colorlet{#1PaleTeal}{UTPPaleTeal}%
  \colorlet{#1PaleOrange}{UTPPaleOrange}\colorlet{#1PaleRed}{UTPPaleRed}%
  \colorlet{#1PaleGray}{UTPPaleGray}%
}
\utpmkalias{ML}
\utpmkalias{ADS}
\utpmkalias{DS}
\utpmkalias{ERP}
\color{UTPInk}

% --------------------------------------------------------------------
%  4.  PAGE LAYOUT AND TABLES
% --------------------------------------------------------------------
%  Paragraph spacing, list indents, and the table look: alternating
%  row tint, top-aligned paragraph columns, ragged-right cells.
%
%  A book typeset tighter overrides \arraystretch, \tabcolsep or
%  \emergencystretch after the input. Data Communication is the one
%  that currently does.

\setlength{\parindent}{0pt}
\setlength{\parskip}{0.55em}
\setlist[itemize]{leftmargin=1.45em,itemsep=0.22em,topsep=0.25em}
\setlist[enumerate]{leftmargin=1.75em,itemsep=0.28em,topsep=0.3em}
\setlist[description]{font=\bfseries\color{UTPNavy},itemsep=0.25em,topsep=0.3em}
\renewcommand{\arraystretch}{1.34}
\setlength{\tabcolsep}{6.5pt}
\arrayrulecolor{UTPBlue!45}
% Keep paragraph-style columns top-aligned and ragged. Tables mix fixed p
% columns with flexible X columns; giving X the same baseline stops one
% column floating halfway down a multi-line row.
\renewcommand{\tabularxcolumn}[1]{p{#1}}
\makeatletter
\g@addto@macro\@arrayparboxrestore{\RaggedRight\arraybackslash}
\makeatother
\AtBeginEnvironment{tabular}{\small\setlength{\extrarowheight}{0.6pt}\rowcolors{2}{UTPPaleBlue!72}{white}}
\AtBeginEnvironment{tabularx}{\small\setlength{\extrarowheight}{0.6pt}\rowcolors{2}{UTPPaleBlue!72}{white}}
\AtBeginEnvironment{figure}{\centering}
\emergencystretch=1.5em
\hbadness=10000
\hfuzz=8pt

\pgfplotsset{
  compat=1.18,
  every axis/.append style={
    axis line style={draw=UTPMuted!85,line width=0.55pt},
    tick style={draw=UTPMuted!85},
    tick label style={font=\scriptsize\color{UTPMuted}},
    label style={font=\small\color{UTPNavy}},
    title style={font=\small\bfseries\color{UTPNavy}},
    legend style={font=\scriptsize,draw=UTPBlue!22,fill=white,rounded corners=1pt}
  },
  every axis plot/.append style={line width=1.15pt}
}
\tikzset{every picture/.append style={line cap=round,line join=round}}

% --------------------------------------------------------------------
%  5.  TABLE OF CONTENTS
% --------------------------------------------------------------------
%  Formatting for part, chapter and lab entries. The lab level is a
%  custom tocloft list so lab banners appear in the contents.
%
%  Part styling is skipped when \utpstyleparts is 0, which is what a
%  book without \part divisions should set.

\addto\captionsamerican{\renewcommand{\contentsname}{Learning Path}}
\newlistentry{lab}{toc}{-1}
\providecommand{\utpstyleparts}{1}
\ifnum\utpstyleparts=1
\renewcommand{\cftpartpresnum}{}
\renewcommand{\cftpartaftersnum}{}
\setlength{\cftpartnumwidth}{5.2em}
\setlength{\cftbeforepartskip}{0.38em}
\renewcommand{\cftpartfont}{\normalsize\bfseries\color{UTPNavy}}
\renewcommand{\cftpartpagefont}{\normalsize\bfseries\color{UTPNavy}}
\renewcommand{\cftpartleader}{\hfill}
\fi
\renewcommand{\cftchappresnum}{Chapter~}
\renewcommand{\cftchapaftersnum}{\quad}
\setlength{\cftchapindent}{1.5em}
\setlength{\cftchapnumwidth}{6em}
\setlength{\cftbeforechapskip}{0pt}
\renewcommand{\cftchapfont}{\normalfont}
\renewcommand{\cftchappagefont}{\normalfont}
\renewcommand{\cftchapleader}{\cftdotfill{\cftdotsep}}
\setlength{\cftlabindent}{3em}
\setlength{\cftlabnumwidth}{0em}
\setlength{\cftbeforelabskip}{0.10em}
\renewcommand{\cftlabfont}{\bfseries\color{UTPOrange!90!black}}
\renewcommand{\cftlabpagefont}{\bfseries\color{UTPOrange!90!black}}
\renewcommand{\cftlableader}{\color{UTPOrange!65!UTPMuted}\cftdotfill{\cftlabdotsep}}

% --------------------------------------------------------------------
%  6.  HEADINGS
% --------------------------------------------------------------------
%  The chapter badge (a coloured block with a word and the number),
%  part pages, and section spacing.
%
%  The word in the badge is \utpchapterlabel; Data Communication sets
%  it to LAB because each of its chapters is a lab.

\titleformat{\chapter}[display]
  {\bfseries\color{UTPNavy}}
  {\filleft\colorbox{UTPBlue}{\parbox{2.55cm}{\centering\color{white}\Large
    \utpchapterlabel\\[-0.15em]\Huge\thechapter}}}
  {0.65em}{\Huge\raggedright}
  [\vspace{0.25em}{\color{UTPBlue}\titlerule[1.2pt]}]
\renewcommand{\thepart}{Stage~\arabic{part}}
\titleformat{\part}[display]
  {\centering\color{UTPNavy}\bfseries}
  {\Large\color{UTPTeal}\thepart}{0.45em}{\Huge}
  [\vspace{0.55em}\color{UTPTeal}\rule{0.45\linewidth}{1.5pt}]
\titleformat{\section}{\Large\bfseries\color{UTPNavy}}{\thesection}{0.55em}{}
\titleformat{\subsection}{\large\bfseries\color{UTPBlue}}{\thesubsection}{0.5em}{}
\titlespacing*{\chapter}{0pt}{-1.8em}{1.0em}
\titlespacing*{\part}{0pt}{0pt}{1em}
\titlespacing*{\section}{0pt}{1.2em}{0.35em}
\titlespacing*{\subsection}{0pt}{0.95em}{0.2em}

% --------------------------------------------------------------------
%  7.  PAGE FURNITURE
% --------------------------------------------------------------------
%  Running heads and footers. The left header is \utprunninghead;
%  the right one follows the current chapter.

\pagestyle{fancy}
\fancyhf{}
\lhead{\small\color{UTPMuted}\utprunninghead}
\rhead{\small\color{UTPMuted}\nouppercase{\leftmark}}
\cfoot{\small\color{UTPMuted}\thepage}
\renewcommand{\headrulewidth}{0.5pt}
\renewcommand{\headrule}{\hbox to\headwidth{\color{UTPPaleBlue}\leaders\hrule height \headrulewidth\hfill}}
\renewcommand{\chaptermark}[1]{\markboth{#1}{}}
\fancypagestyle{plain}{%
  \fancyhf{}%
  \cfoot{\small\color{UTPMuted}\thepage}%
  \renewcommand{\headrulewidth}{0pt}%
}

% --------------------------------------------------------------------
%  8.  LEARNING BOXES
% --------------------------------------------------------------------
%  The coloured callouts the chapters use. Every title is a macro, so
%  a book can rename a box without restating its geometry:
%
%      \def\utptitlepitfall{Technical note}   % before the input
%
%  To change how a box LOOKS in one book, use \renewtcolorbox after
%  the input instead. To add a box only one book needs, declare it
%  there with \newtcolorbox.

% Box titles -- override any of these before the input line.
\providecommand{\utptitlekeyidea}{Key idea}
\providecommand{\utptitleworked}{Worked example}
\providecommand{\utptitlepitfall}{Common mistake}
\providecommand{\utptitletryit}{Try it yourself}
\providecommand{\utptitlequickcheck}{Check your understanding}
\providecommand{\utptitlerecap}{Chapter recap}
\providecommand{\utptitledeliverable}{Lab deliverable}
\providecommand{\utptitlelabmode}{How much is scaffolded?}
\providecommand{\utptitleintuition}{Plain idea}
\providecommand{\utptitlefirstread}{Core path}
\providecommand{\utptitleglance}{Chapter at a glance}
\providecommand{\utptitledebug}{Debugging task}

% Explaining an idea
\newtcolorbox{keyidea}{enhanced,breakable,colback=UTPPaleTeal,colframe=UTPTeal, boxrule=0.8pt,arc=2mm,left=2.5mm,right=2.5mm,top=1.8mm,bottom=1.8mm, fonttitle=\bfseries\color{UTPNavy},title=\utptitlekeyidea}
\newtcolorbox{workedexample}{enhanced,breakable,colback=UTPPaleBlue,colframe=UTPBlue, boxrule=0.8pt,arc=2mm,left=2.5mm,right=2.5mm,top=1.8mm,bottom=1.8mm, fonttitle=\bfseries\color{UTPNavy},title=\utptitleworked}
\newtcolorbox{pitfall}{enhanced,breakable,colback=UTPPaleRed,colframe=UTPRed, boxrule=0.8pt,arc=2mm,left=2.5mm,right=2.5mm,top=1.8mm,bottom=1.8mm, fonttitle=\bfseries\color{UTPNavy},title=\utptitlepitfall}
% Asking the reader to do something
\newtcolorbox{tryit}{enhanced,breakable,colback=white,colframe=UTPTeal, boxrule=0.9pt,arc=2mm,left=2.5mm,right=2.5mm,top=1.8mm,bottom=1.8mm, fonttitle=\bfseries\color{white},colbacktitle=UTPTeal,title=\utptitletryit}
\newtcolorbox{quickcheck}{enhanced,breakable,colback=white,colframe=UTPOrange, boxrule=0.9pt,arc=2mm,left=2.5mm,right=2.5mm,top=1.8mm,bottom=1.8mm, fonttitle=\bfseries\color{white},colbacktitle=UTPOrange,title=\utptitlequickcheck}
\newtcolorbox{chapterrecap}{enhanced,breakable,colback=UTPPaleTeal,colframe=UTPGreen, boxrule=0.9pt,arc=2mm,left=2.5mm,right=2.5mm,top=1.8mm,bottom=1.8mm, fonttitle=\bfseries\color{white},colbacktitle=UTPGreen,title=\utptitlerecap}
% Lab structure
\newtcolorbox{labdeliverable}{enhanced,breakable,colback=white,colframe=UTPGreen, boxrule=0.9pt,arc=2mm,left=2.6mm,right=2.6mm,top=1.9mm,bottom=1.9mm, fonttitle=\bfseries\color{white},colbacktitle=UTPGreen,title=\utptitledeliverable}
\newtcolorbox{labmode}{enhanced,breakable,colback=UTPPaleBlue!52,colframe=UTPBlue, boxrule=0.75pt,arc=2mm,left=2.6mm,right=2.6mm,top=1.7mm,bottom=1.7mm, fonttitle=\bfseries\color{white},colbacktitle=UTPBlue,title=\utptitlelabmode}
% Softer, lighter-weight asides
\newtcolorbox{intuition}{enhanced,breakable,colback=white,colframe=UTPOrange!42, boxrule=0.5pt,arc=2mm,left=2.5mm,right=2.5mm,top=1.5mm,bottom=1.5mm, fonttitle=\bfseries\color{UTPOrange!85!black},title=\utptitleintuition}
\newtcolorbox{firstread}{enhanced,breakable,colback=white,colframe=UTPBlue!42, boxrule=0.5pt,arc=2mm,left=2.5mm,right=2.5mm,top=1.5mm,bottom=1.5mm, fonttitle=\bfseries\color{UTPBlue}, title=\utptitlefirstread}
\newtcolorbox{debugtask}{enhanced,breakable,colback=UTPPaleRed!58,colframe=UTPRed, boxrule=0.85pt,arc=2mm,left=2.6mm,right=2.6mm,top=1.9mm,bottom=1.9mm, fonttitle=\bfseries\color{white},colbacktitle=UTPRed,title=\utptitledebug}
\newtcolorbox{analogy}{enhanced,breakable,colback=white,colframe=UTPOrange!38, boxrule=0.5pt,arc=2mm,left=2.5mm,right=2.5mm,top=1.5mm,bottom=1.5mm, fonttitle=\bfseries\color{UTPOrange!85!black},title=Analogy}
\newtcolorbox{coursechoice}{enhanced,breakable,colback=white,colframe=UTPNavy, boxrule=0.8pt,arc=2mm,left=2.8mm,right=2.8mm,top=2mm,bottom=2mm, fonttitle=\bfseries\color{white},colbacktitle=UTPNavy,title=Course choice}
\newtcolorbox{decisionmemo}{enhanced,breakable,colback=UTPPaleOrange!55,colframe=UTPOrange, boxrule=0.85pt,arc=2mm,left=2.6mm,right=2.6mm,top=1.9mm,bottom=1.9mm, fonttitle=\bfseries\color{white},colbacktitle=UTPOrange,title=Decision memo}
\newtcolorbox{libnote}{enhanced,breakable,colback=white,colframe=UTPMuted!48, boxrule=0.5pt,arc=2mm,left=2.5mm,right=2.5mm,top=1.5mm,bottom=1.5mm, fonttitle=\bfseries\color{UTPNavy},title=Reference note}
\newtcolorbox{miniproject}{enhanced,breakable,colback=UTPPaleTeal!52,colframe=UTPGreen, boxrule=0.9pt,arc=2mm,left=2.6mm,right=2.6mm,top=1.9mm,bottom=1.9mm, fonttitle=\bfseries\color{white},colbacktitle=UTPGreen,title=Mini-project}
% Boxes that take a title argument
\newtcolorbox{modelcard}[1]{enhanced,breakable,colback=white,colframe=UTPBlue, boxrule=0.9pt,arc=2mm,left=2.5mm,right=2.5mm,top=1.8mm,bottom=1.8mm, fonttitle=\bfseries\color{white},colbacktitle=UTPBlue,title=#1}
\newtcolorbox{optionaltopic}{enhanced,breakable,colback=white,colframe=UTPMuted!58, boxrule=0.5pt,arc=2mm,left=2.5mm,right=2.5mm,top=1.5mm,bottom=1.5mm, fonttitle=\bfseries\color{UTPMuted},title=Beyond the core}
\newtcolorbox{secondpass}{enhanced,breakable,colback=UTPPaleGray!82,colframe=UTPMuted, boxrule=0.75pt,arc=2mm,left=2.6mm,right=2.6mm,top=1.9mm,bottom=1.9mm, fonttitle=\bfseries\color{white},colbacktitle=UTPMuted,title=Second pass}
\newtcolorbox{labproject}[1]{enhanced,breakable,colback=white,colframe=UTPBlue, boxrule=0.9pt,arc=2mm,left=2.6mm,right=2.6mm,top=1.9mm,bottom=1.9mm, fonttitle=\bfseries\color{white},colbacktitle=UTPBlue,title={#1}}
% Frames for a figure or a table
\newtcolorbox{mlfigure}{enhanced,width=\linewidth,colback=white,colframe=UTPBlue!28, boxrule=0.7pt,arc=2mm, left=2.2mm,right=2.2mm,top=2.2mm,bottom=1.8mm, before skip=0.75em,after skip=0.85em, before upper={\centering}, halign=center}
\newtcolorbox{mltable}{enhanced,width=\linewidth,colback=white,colframe=UTPTeal!32, boxrule=0.7pt,arc=2mm, left=2.2mm,right=2.2mm,top=1.8mm,bottom=1.8mm, before skip=0.75em,after skip=0.85em, halign=center}

% --------------------------------------------------------------------
%  9.  SMALL COMMANDS
% --------------------------------------------------------------------
%  \term for a defined term, \cmd for inline code, \chapref and
%  \secref for cross-references, and the mlsteps numbered-step list.
%
%      \term{overfitting}          a term being defined
%      \cmd{git status}            inline code
%      \chapref{ch:trees}          "Chapter 12"
%      \secref{sec:pruning}        "Section 12.3"
%      \chapterguide{prior}{outcomes}   the box that opens a chapter
%      \begin{mlsteps} ... \end{mlsteps}  a numbered Step 1/Step 2 list

\newcommand{\term}[1]{\textbf{\color{UTPNavy}#1}}
\newcommand{\cmd}[1]{\texttt{#1}}
\newcommand{\chapref}[1]{Chapter~\ref{#1}}
\newcommand{\secref}[1]{Section~\ref{#1}}
\newcommand{\chapterguide}[2]{%
  \begin{tcolorbox}[enhanced,breakable,colback=white,colframe=UTPNavy,boxrule=0.8pt,
    arc=2mm,left=2.5mm,right=2.5mm,top=2mm,bottom=2mm,
    title=\utptitleglance,fonttitle=\bfseries\color{white},colbacktitle=UTPNavy]
    \textbf{Before you start:} #1\par
    \textbf{By the end, you can:} #2
  \end{tcolorbox}}
\newenvironment{mlsteps}
  {\begin{enumerate}[label=\textcolor{UTPBlue}{\bfseries Step~\arabic*:},leftmargin=4.5em,itemsep=0.38em]}
  {\end{enumerate}}

% --------------------------------------------------------------------
%  10.  LAB MILESTONE BANNER
% --------------------------------------------------------------------
%  The full-width page break that opens each lab. Called from the lab
%  files as \labsection{number}{title}{label}.
%
%  The "of N" comes from \utplabtotal, so a book states its lab count
%  once in its configuration rather than in every lab file.

\newcommand{\labsection}[3]{%
  \clearpage
  \phantomsection
  \addcontentsline{toc}{lab}{Lab #1 of \utplabtotal\space--- #2}%
  \markboth{Lab #1 of \utplabtotal\space--- #2}{}%
  \begin{tcolorbox}[enhanced,width=\linewidth,colback=UTPNavy,colframe=UTPNavy,
    boxrule=0pt,arc=2mm,left=5mm,right=5mm,top=5mm,bottom=5mm,
    before skip=0pt,after skip=1.2em]
    {\large\bfseries\color{UTPTeal!35!white} PRACTICAL MILESTONE\par}
    \vspace{0.25em}
    {\Huge\bfseries\color{white} Lab #1 of \utplabtotal\par}
    \vspace{0.20em}
    {\Large\bfseries\color{white} #2\par}
  \end{tcolorbox}
  \label{#3}%
}

% --------------------------------------------------------------------
%  11.  STAGE DIVIDER
% --------------------------------------------------------------------
%  A full-page divider used by the longer books to open a part:
%  \substage{3A}{Supervised Workflow}{one-line summary}.
%  Harmless to leave unused.

\newcommand{\substage}[3]{%
  \cleardoublepage
  \thispagestyle{empty}
  \phantomsection
  \addcontentsline{toc}{chapter}{Stage #1 --- #2}
  \begin{center}
    \vspace*{0.20\textheight}
    {\Large\bfseries\color{UTPTeal} Stage #1\par}
    \vspace{0.45em}
    {\Huge\bfseries\color{UTPNavy} #2\par}
    \vspace{0.65em}
    {\color{UTPTeal}\rule{0.46\linewidth}{1.5pt}\par}
    \vspace{1.1em}
    \begin{minipage}{0.76\linewidth}
      \centering\large\color{UTPMuted} #3
    \end{minipage}
  \end{center}
  \cleardoublepage
}

% --------------------------------------------------------------------
%  12.  LISTING STYLES
% --------------------------------------------------------------------
%  Two base styles. A course builds its own on top of them so the
%  colours and frame stay consistent:
%
%      \lstdefinestyle{mlpython}{style=utpcode,language=Python}
%      \lstdefinestyle{mlterminal}{style=utpterminal}

\lstdefinestyle{utpcode}{
  basicstyle=\ttfamily\scriptsize,
  keywordstyle=\bfseries\color{UTPBlue},
  commentstyle=\itshape\color{UTPMuted},
  stringstyle=\color{UTPGreen!70!black},
  backgroundcolor=\color{UTPPaleGray!72},
  frame=single,
  rulecolor=\color{UTPBlue!28},
  framesep=5pt,
  breaklines=true,
  breakatwhitespace=false,
  showstringspaces=false,
  keepspaces=true,
  columns=fullflexible,
  upquote=true,
  tabsize=4,
  aboveskip=0.7em,
  belowskip=0.8em
}
\lstdefinestyle{utpterminal}{
  basicstyle=\ttfamily\scriptsize,
  backgroundcolor=\color{UTPPaleGray!72},
  frame=single,
  rulecolor=\color{UTPMuted!35},
  framesep=5pt,
  breaklines=true,
  showstringspaces=false,
  keepspaces=true,
  columns=fullflexible,
  upquote=true,
  aboveskip=0.6em,
  belowskip=0.7em
}

\lstdefinestyle{adscpp}{
  language=C++,
  basicstyle=\ttfamily\scriptsize,
  keywordstyle=\bfseries\color{UTPBlue},
  commentstyle=\itshape\color{UTPMuted},
  stringstyle=\color{UTPGreen!70!black},
  backgroundcolor=\color{UTPPaleGray!72},
  frame=single,
  rulecolor=\color{UTPBlue!28},
  framesep=5pt,
  breaklines=true,
  breakatwhitespace=false,
  showstringspaces=false,
  keepspaces=true,
  columns=fullflexible,
  upquote=true,
  tabsize=4,
  aboveskip=0.7em,
  belowskip=0.8em
}

\lstdefinestyle{mlpython}{
  language=Python,
  basicstyle=\ttfamily\scriptsize,
  keywordstyle=\bfseries\color{UTPBlue},
  commentstyle=\itshape\color{UTPMuted},
  stringstyle=\color{UTPGreen!70!black},
  backgroundcolor=\color{UTPPaleGray!72},
  frame=single,
  rulecolor=\color{UTPBlue!28},
  framesep=5pt,
  breaklines=true,
  breakatwhitespace=false,
  showstringspaces=false,
  keepspaces=true,
  columns=fullflexible,
  upquote=true,
  tabsize=4,
  aboveskip=0.7em,
  belowskip=0.8em
}

\lstdefinestyle{mlmatlab}{
  language=Matlab,
  basicstyle=\ttfamily\scriptsize,
  keywordstyle=\bfseries\color{UTPBlue},
  commentstyle=\itshape\color{UTPMuted},
  stringstyle=\color{UTPGreen!70!black},
  backgroundcolor=\color{UTPPaleGray!72},
  frame=single,
  rulecolor=\color{UTPBlue!28},
  framesep=5pt,
  breaklines=true,
  breakatwhitespace=false,
  showstringspaces=false,
  keepspaces=true,
  columns=fullflexible,
  upquote=true,
  tabsize=4,
  aboveskip=0.7em,
  belowskip=0.8em
}

\lstdefinestyle{mlterminal}{
  basicstyle=\ttfamily\scriptsize,
  backgroundcolor=\color{UTPPaleGray!72},
  frame=single,
  rulecolor=\color{UTPMuted!35},
  framesep=5pt,
  breaklines=true,
  showstringspaces=false,
  keepspaces=true,
  columns=fullflexible,
  upquote=true,
  aboveskip=0.6em,
  belowskip=0.7em
}

  \lstdefinestyle{dcnterminal}{
  basicstyle=\ttfamily\scriptsize,
  backgroundcolor=\color{UTPPaleGray!72},
  frame=single,rulecolor=\color{UTPMuted!35},
  framesep=5pt,
  breaklines=true,
  showstringspaces=false,
  keepspaces=true,
  columns=fullflexible,
  upquote=true,
  aboveskip=0.6em,
  belowskip=0.7em
}

\lstdefinestyle{dsr}{
  language=R,
  basicstyle=\ttfamily\scriptsize,
  keywordstyle=\bfseries\color{UTPBlue},
  commentstyle=\itshape\color{UTPMuted},
  stringstyle=\color{UTPGreen!70!black},
  backgroundcolor=\color{UTPPaleGray!72},
  frame=single,
  rulecolor=\color{UTPBlue!28},
  framesep=5pt,
  breaklines=true,
  breakatwhitespace=false,
  showstringspaces=false,
  keepspaces=true,
  columns=fullflexible,
  upquote=true,
  tabsize=2,
  aboveskip=0.7em,
  belowskip=0.8em
}

\lstdefinestyle{dsterminal}{
  basicstyle=\ttfamily\scriptsize,
  backgroundcolor=\color{UTPPaleGray!72},
  frame=single,
  rulecolor=\color{UTPMuted!35},
  framesep=5pt,
  breaklines=true,
  showstringspaces=false,
  keepspaces=true,
  columns=fullflexible,
  upquote=true,
  aboveskip=0.6em,
  belowskip=0.7em
}

\lstdefinestyle{erpterminal}{
  basicstyle=\ttfamily\scriptsize,
  backgroundcolor=\color{UTPPaleGray!72},
  frame=single,
  rulecolor=\color{UTPBlue!28},
  framesep=5pt,
  breaklines=true,
  showstringspaces=false,
  keepspaces=true,
  columns=fullflexible,
  upquote=true,
  tabsize=4,
  aboveskip=0.7em,
  belowskip=0.8em
}

\lstdefinestyle{mlpython}{
  language=Python,
  basicstyle=\ttfamily\scriptsize,
  keywordstyle=\bfseries\color{UTPBlue},
  commentstyle=\itshape\color{UTPMuted},
  stringstyle=\color{UTPGreen!70!black},
  backgroundcolor=\color{UTPPaleGray!72},
  frame=single,
  rulecolor=\color{UTPBlue!28},
  framesep=5pt,
  breaklines=true,
  breakatwhitespace=false,
  showstringspaces=false,
  keepspaces=true,
  columns=fullflexible,
  upquote=true,
  tabsize=4,
  aboveskip=0.7em,
  belowskip=0.8em
}
\lstdefinestyle{mlterminal}{
  basicstyle=\ttfamily\scriptsize,
  backgroundcolor=\color{UTPPaleGray!72},
  frame=single,
  rulecolor=\color{UTPMuted!35},
  framesep=5pt,
  breaklines=true,
  showstringspaces=false,
  keepspaces=true,
  columns=fullflexible,
  upquote=true,
  aboveskip=0.6em,
  belowskip=0.7em
}

\lstdefinestyle{mlpython}{
  language=Python,
  basicstyle=\ttfamily\scriptsize,
  keywordstyle=\bfseries\color{UTPBlue},
  commentstyle=\itshape\color{UTPMuted},
  stringstyle=\color{UTPGreen!70!black},
  backgroundcolor=\color{UTPPaleGray!72},
  frame=single,
  rulecolor=\color{UTPBlue!28},
  framesep=5pt,
  breaklines=true,
  breakatwhitespace=false,
  showstringspaces=false,
  keepspaces=true,
  columns=fullflexible,
  upquote=true,
  tabsize=4,
  aboveskip=0.7em,
  belowskip=0.8em
}
\lstdefinestyle{mlterminal}{
  basicstyle=\ttfamily\scriptsize,
  backgroundcolor=\color{UTPPaleGray!72},
  frame=single,
  rulecolor=\color{UTPMuted!35},
  framesep=5pt,
  breaklines=true,
  showstringspaces=false,
  keepspaces=true,
  columns=fullflexible,
  upquote=true,
  aboveskip=0.6em,
  belowskip=0.7em
}
\lstdefinestyle{osc}{
  language=C,
  basicstyle=\ttfamily\scriptsize,
  keywordstyle=\bfseries\color{UTPBlue},
  commentstyle=\itshape\color{UTPMuted},
  stringstyle=\color{UTPGreen!70!black},
  backgroundcolor=\color{UTPPaleGray!72},
  frame=single,
  rulecolor=\color{UTPBlue!28},
  framesep=5pt,
  breaklines=true,
  showstringspaces=false,
  keepspaces=true,
  columns=fullflexible,
  upquote=true,
  tabsize=4,
  aboveskip=0.7em,
  belowskip=0.8em
}
\lstdefinestyle{osbash}{
  language=bash,
  basicstyle=\ttfamily\scriptsize,
  keywordstyle=\bfseries\color{UTPBlue},
  commentstyle=\itshape\color{UTPMuted},
  stringstyle=\color{UTPGreen!70!black},
  backgroundcolor=\color{UTPPaleGray!72},
  frame=single,
  rulecolor=\color{UTPMuted!35},
  framesep=5pt,
  breaklines=true,
  showstringspaces=false,
  keepspaces=true,
  columns=fullflexible,
  upquote=true,
  aboveskip=0.6em,
  belowskip=0.7em
}



%      \graphicspath{{../.shared/}}
%
%  Configuration goes BEFORE the \input; overrides go AFTER it.
%  Paths are relative, so compile from inside the course folder.
%
%  Adding a new course? Copy .shared/TEMPLATE-main.tex and follow
%  .shared/ADDING-A-COURSE.md. Every step is a numbered checklist.
%
%  ------------------------------------------------------------------
%  CONTENTS
%  ------------------------------------------------------------------
%    1  Course configuration     the knobs a book may set
%    2  Packages                 the shared package set
%    3  Colour palette           one palette, plus per-course aliases
%    4  Page layout and tables   margins, lists, table styling
%    5  Table of contents        entry formatting for parts/chapters/labs
%    6  Headings                 chapter badge, part pages, sections
%    7  Page furniture           running heads and footers
%    8  Learning boxes           key idea, pitfall, try it, and friends
%    9  Small commands           \term, \cmd, \chapterguide, \mlsteps
%   10  Lab milestone banner     the full-width "Lab N of M" page break
%   11  Stage divider            the part-opening page in longer books
%   12  Listing styles           the base code and terminal styles
% =====================================================================

% --------------------------------------------------------------------
%  1.  COURSE CONFIGURATION
% --------------------------------------------------------------------
%  The three knobs below, plus \utpstyleparts in section 5 and the box
%  titles in section 8, are the whole configuration surface. A book
%  sets what it needs BEFORE the \input and inherits the rest.
%
%    \utprunninghead   the left running head on every page.
%    \utpchapterlabel  the word in the chapter badge. CHAPTER by
%                      default; a book whose chapters are labs sets LAB.
%    \utplabtotal      the N in the "Lab 3 of N" banners.
%
%  \providecommand means "define unless the book already did", which is
%  what makes the \def-before-\input pattern work.

\providecommand{\utprunninghead}{Course Notes}
\providecommand{\utpchapterlabel}{CHAPTER}
\providecommand{\utplabtotal}{10}

% --------------------------------------------------------------------
%  2.  PACKAGES
% --------------------------------------------------------------------
%  The union of what the eight books need. A course that needs
%  something extra loads it in its own file, under the heading
%  "packages only this book needs" (algorithm2e, caption, float and
%  multirow are the only current examples).
%
%  The tikz library list is the union across all books, so a diagram
%  copied from one course compiles in another.

\usepackage[T1]{fontenc}
\IfFileExists{glyphtounicode.tex}{\input{glyphtounicode}\pdfgentounicode=1}{}
\usepackage[american]{babel}
\IfFileExists{sourcesanspro.sty}
  {\usepackage[default]{sourcesanspro}}
  {\usepackage[scaled]{helvet}\renewcommand{\familydefault}{\sfdefault}}
\DeclareFontShape{T1}{cmss}{b}{n}{<->ssub*cmss/bx/n}{}
\DeclareFontShape{TS1}{cmss}{b}{n}{<->ssub*cmss/bx/n}{}
\usepackage[margin=21mm,headheight=22pt]{geometry}
\usepackage{amsmath,amssymb,mathtools,bm}
\usepackage{microtype}
\usepackage{enumitem}
\usepackage{booktabs}
\usepackage{tabularx}
\usepackage{array}
\usepackage{ragged2e}
\usepackage{etoolbox}
\usepackage{titlesec}
\usepackage{tocloft}
\usepackage{fancyhdr}
\usepackage[table]{xcolor}
\usepackage{tcolorbox}
\tcbuselibrary{breakable,skins}
\usepackage[hidelinks]{hyperref}
\usepackage{bookmark}
\usepackage{listings}
\usepackage{graphicx}
\usepackage{tikz}
\usetikzlibrary{angles,arrows.meta,backgrounds,calc,chains,decorations.pathmorphing,decorations.pathreplacing,fit,matrix,patterns,positioning,quotes,shadows.blur,shapes.geometric,shapes.misc}
\usepackage{pgfplots}

% --------------------------------------------------------------------
%  3.  COLOUR PALETTE
% --------------------------------------------------------------------
%  One palette, defined once under canonical UTP* names. Use these
%  in any new material.
%
%  The existing chapters were written against course-prefixed names
%  (MLNavy, ADSPaleBlue, ...) in over two thousand places, so rather
%  than rewrite them the block below aliases each prefix onto the
%  canonical colours. All four sets are always defined, so a chapter
%  moved between courses still resolves.

\definecolor{UTPNavy}{HTML}{18324A}
\definecolor{UTPBlue}{HTML}{246B9E}
\definecolor{UTPTeal}{HTML}{138A80}
\definecolor{UTPGreen}{HTML}{2D7D46}
\definecolor{UTPOrange}{HTML}{C86B24}
\definecolor{UTPRed}{HTML}{B34843}
\definecolor{UTPGold}{HTML}{D4A13E}
\definecolor{UTPInk}{HTML}{263238}
\definecolor{UTPMuted}{HTML}{60717C}
\definecolor{UTPPaleBlue}{HTML}{EEF6FB}
\definecolor{UTPPaleTeal}{HTML}{EEF8F6}
\definecolor{UTPPaleOrange}{HTML}{FFF5EA}
\definecolor{UTPPaleRed}{HTML}{FFF0EF}
\definecolor{UTPPaleGray}{HTML}{F4F6F7}

% Course files were written against prefixed colour names. Rather than edit
% 2000+ references across the chapters, alias each prefix onto the canonical
% palette. Only these four prefixes appear anywhere in the books.
\def\utpmkalias#1{%
  \colorlet{#1Navy}{UTPNavy}\colorlet{#1Blue}{UTPBlue}%
  \colorlet{#1Teal}{UTPTeal}\colorlet{#1Green}{UTPGreen}%
  \colorlet{#1Orange}{UTPOrange}\colorlet{#1Red}{UTPRed}%
  \colorlet{#1Ink}{UTPInk}\colorlet{#1Muted}{UTPMuted}%
  \colorlet{#1PaleBlue}{UTPPaleBlue}\colorlet{#1PaleTeal}{UTPPaleTeal}%
  \colorlet{#1PaleOrange}{UTPPaleOrange}\colorlet{#1PaleRed}{UTPPaleRed}%
  \colorlet{#1PaleGray}{UTPPaleGray}%
}
\utpmkalias{ML}
\utpmkalias{ADS}
\utpmkalias{DS}
\utpmkalias{ERP}
\color{UTPInk}

% --------------------------------------------------------------------
%  4.  PAGE LAYOUT AND TABLES
% --------------------------------------------------------------------
%  Paragraph spacing, list indents, and the table look: alternating
%  row tint, top-aligned paragraph columns, ragged-right cells.
%
%  A book typeset tighter overrides \arraystretch, \tabcolsep or
%  \emergencystretch after the input. Data Communication is the one
%  that currently does.

\setlength{\parindent}{0pt}
\setlength{\parskip}{0.55em}
\setlist[itemize]{leftmargin=1.45em,itemsep=0.22em,topsep=0.25em}
\setlist[enumerate]{leftmargin=1.75em,itemsep=0.28em,topsep=0.3em}
\setlist[description]{font=\bfseries\color{UTPNavy},itemsep=0.25em,topsep=0.3em}
\renewcommand{\arraystretch}{1.34}
\setlength{\tabcolsep}{6.5pt}
\arrayrulecolor{UTPBlue!45}
% Keep paragraph-style columns top-aligned and ragged. Tables mix fixed p
% columns with flexible X columns; giving X the same baseline stops one
% column floating halfway down a multi-line row.
\renewcommand{\tabularxcolumn}[1]{p{#1}}
\makeatletter
\g@addto@macro\@arrayparboxrestore{\RaggedRight\arraybackslash}
\makeatother
\AtBeginEnvironment{tabular}{\small\setlength{\extrarowheight}{0.6pt}\rowcolors{2}{UTPPaleBlue!72}{white}}
\AtBeginEnvironment{tabularx}{\small\setlength{\extrarowheight}{0.6pt}\rowcolors{2}{UTPPaleBlue!72}{white}}
\AtBeginEnvironment{figure}{\centering}
\emergencystretch=1.5em
\hbadness=10000
\hfuzz=8pt

\pgfplotsset{
  compat=1.18,
  every axis/.append style={
    axis line style={draw=UTPMuted!85,line width=0.55pt},
    tick style={draw=UTPMuted!85},
    tick label style={font=\scriptsize\color{UTPMuted}},
    label style={font=\small\color{UTPNavy}},
    title style={font=\small\bfseries\color{UTPNavy}},
    legend style={font=\scriptsize,draw=UTPBlue!22,fill=white,rounded corners=1pt}
  },
  every axis plot/.append style={line width=1.15pt}
}
\tikzset{every picture/.append style={line cap=round,line join=round}}

% --------------------------------------------------------------------
%  5.  TABLE OF CONTENTS
% --------------------------------------------------------------------
%  Formatting for part, chapter and lab entries. The lab level is a
%  custom tocloft list so lab banners appear in the contents.
%
%  Part styling is skipped when \utpstyleparts is 0, which is what a
%  book without \part divisions should set.

\addto\captionsamerican{\renewcommand{\contentsname}{Learning Path}}
\newlistentry{lab}{toc}{-1}
\providecommand{\utpstyleparts}{1}
\ifnum\utpstyleparts=1
\renewcommand{\cftpartpresnum}{}
\renewcommand{\cftpartaftersnum}{}
\setlength{\cftpartnumwidth}{5.2em}
\setlength{\cftbeforepartskip}{0.38em}
\renewcommand{\cftpartfont}{\normalsize\bfseries\color{UTPNavy}}
\renewcommand{\cftpartpagefont}{\normalsize\bfseries\color{UTPNavy}}
\renewcommand{\cftpartleader}{\hfill}
\fi
\renewcommand{\cftchappresnum}{Chapter~}
\renewcommand{\cftchapaftersnum}{\quad}
\setlength{\cftchapindent}{1.5em}
\setlength{\cftchapnumwidth}{6em}
\setlength{\cftbeforechapskip}{0pt}
\renewcommand{\cftchapfont}{\normalfont}
\renewcommand{\cftchappagefont}{\normalfont}
\renewcommand{\cftchapleader}{\cftdotfill{\cftdotsep}}
\setlength{\cftlabindent}{3em}
\setlength{\cftlabnumwidth}{0em}
\setlength{\cftbeforelabskip}{0.10em}
\renewcommand{\cftlabfont}{\bfseries\color{UTPOrange!90!black}}
\renewcommand{\cftlabpagefont}{\bfseries\color{UTPOrange!90!black}}
\renewcommand{\cftlableader}{\color{UTPOrange!65!UTPMuted}\cftdotfill{\cftlabdotsep}}

% --------------------------------------------------------------------
%  6.  HEADINGS
% --------------------------------------------------------------------
%  The chapter badge (a coloured block with a word and the number),
%  part pages, and section spacing.
%
%  The word in the badge is \utpchapterlabel; Data Communication sets
%  it to LAB because each of its chapters is a lab.

\titleformat{\chapter}[display]
  {\bfseries\color{UTPNavy}}
  {\filleft\colorbox{UTPBlue}{\parbox{2.55cm}{\centering\color{white}\Large
    \utpchapterlabel\\[-0.15em]\Huge\thechapter}}}
  {0.65em}{\Huge\raggedright}
  [\vspace{0.25em}{\color{UTPBlue}\titlerule[1.2pt]}]
\renewcommand{\thepart}{Stage~\arabic{part}}
\titleformat{\part}[display]
  {\centering\color{UTPNavy}\bfseries}
  {\Large\color{UTPTeal}\thepart}{0.45em}{\Huge}
  [\vspace{0.55em}\color{UTPTeal}\rule{0.45\linewidth}{1.5pt}]
\titleformat{\section}{\Large\bfseries\color{UTPNavy}}{\thesection}{0.55em}{}
\titleformat{\subsection}{\large\bfseries\color{UTPBlue}}{\thesubsection}{0.5em}{}
\titlespacing*{\chapter}{0pt}{-1.8em}{1.0em}
\titlespacing*{\part}{0pt}{0pt}{1em}
\titlespacing*{\section}{0pt}{1.2em}{0.35em}
\titlespacing*{\subsection}{0pt}{0.95em}{0.2em}

% --------------------------------------------------------------------
%  7.  PAGE FURNITURE
% --------------------------------------------------------------------
%  Running heads and footers. The left header is \utprunninghead;
%  the right one follows the current chapter.

\pagestyle{fancy}
\fancyhf{}
\lhead{\small\color{UTPMuted}\utprunninghead}
\rhead{\small\color{UTPMuted}\nouppercase{\leftmark}}
\cfoot{\small\color{UTPMuted}\thepage}
\renewcommand{\headrulewidth}{0.5pt}
\renewcommand{\headrule}{\hbox to\headwidth{\color{UTPPaleBlue}\leaders\hrule height \headrulewidth\hfill}}
\renewcommand{\chaptermark}[1]{\markboth{#1}{}}
\fancypagestyle{plain}{%
  \fancyhf{}%
  \cfoot{\small\color{UTPMuted}\thepage}%
  \renewcommand{\headrulewidth}{0pt}%
}

% --------------------------------------------------------------------
%  8.  LEARNING BOXES
% --------------------------------------------------------------------
%  The coloured callouts the chapters use. Every title is a macro, so
%  a book can rename a box without restating its geometry:
%
%      \def\utptitlepitfall{Technical note}   % before the input
%
%  To change how a box LOOKS in one book, use \renewtcolorbox after
%  the input instead. To add a box only one book needs, declare it
%  there with \newtcolorbox.

% Box titles -- override any of these before the input line.
\providecommand{\utptitlekeyidea}{Key idea}
\providecommand{\utptitleworked}{Worked example}
\providecommand{\utptitlepitfall}{Common mistake}
\providecommand{\utptitletryit}{Try it yourself}
\providecommand{\utptitlequickcheck}{Check your understanding}
\providecommand{\utptitlerecap}{Chapter recap}
\providecommand{\utptitledeliverable}{Lab deliverable}
\providecommand{\utptitlelabmode}{How much is scaffolded?}
\providecommand{\utptitleintuition}{Plain idea}
\providecommand{\utptitlefirstread}{Core path}
\providecommand{\utptitleglance}{Chapter at a glance}
\providecommand{\utptitledebug}{Debugging task}

% Explaining an idea
\newtcolorbox{keyidea}{enhanced,breakable,colback=UTPPaleTeal,colframe=UTPTeal, boxrule=0.8pt,arc=2mm,left=2.5mm,right=2.5mm,top=1.8mm,bottom=1.8mm, fonttitle=\bfseries\color{UTPNavy},title=\utptitlekeyidea}
\newtcolorbox{workedexample}{enhanced,breakable,colback=UTPPaleBlue,colframe=UTPBlue, boxrule=0.8pt,arc=2mm,left=2.5mm,right=2.5mm,top=1.8mm,bottom=1.8mm, fonttitle=\bfseries\color{UTPNavy},title=\utptitleworked}
\newtcolorbox{pitfall}{enhanced,breakable,colback=UTPPaleRed,colframe=UTPRed, boxrule=0.8pt,arc=2mm,left=2.5mm,right=2.5mm,top=1.8mm,bottom=1.8mm, fonttitle=\bfseries\color{UTPNavy},title=\utptitlepitfall}
% Asking the reader to do something
\newtcolorbox{tryit}{enhanced,breakable,colback=white,colframe=UTPTeal, boxrule=0.9pt,arc=2mm,left=2.5mm,right=2.5mm,top=1.8mm,bottom=1.8mm, fonttitle=\bfseries\color{white},colbacktitle=UTPTeal,title=\utptitletryit}
\newtcolorbox{quickcheck}{enhanced,breakable,colback=white,colframe=UTPOrange, boxrule=0.9pt,arc=2mm,left=2.5mm,right=2.5mm,top=1.8mm,bottom=1.8mm, fonttitle=\bfseries\color{white},colbacktitle=UTPOrange,title=\utptitlequickcheck}
\newtcolorbox{chapterrecap}{enhanced,breakable,colback=UTPPaleTeal,colframe=UTPGreen, boxrule=0.9pt,arc=2mm,left=2.5mm,right=2.5mm,top=1.8mm,bottom=1.8mm, fonttitle=\bfseries\color{white},colbacktitle=UTPGreen,title=\utptitlerecap}
% Lab structure
\newtcolorbox{labdeliverable}{enhanced,breakable,colback=white,colframe=UTPGreen, boxrule=0.9pt,arc=2mm,left=2.6mm,right=2.6mm,top=1.9mm,bottom=1.9mm, fonttitle=\bfseries\color{white},colbacktitle=UTPGreen,title=\utptitledeliverable}
\newtcolorbox{labmode}{enhanced,breakable,colback=UTPPaleBlue!52,colframe=UTPBlue, boxrule=0.75pt,arc=2mm,left=2.6mm,right=2.6mm,top=1.7mm,bottom=1.7mm, fonttitle=\bfseries\color{white},colbacktitle=UTPBlue,title=\utptitlelabmode}
% Softer, lighter-weight asides
\newtcolorbox{intuition}{enhanced,breakable,colback=white,colframe=UTPOrange!42, boxrule=0.5pt,arc=2mm,left=2.5mm,right=2.5mm,top=1.5mm,bottom=1.5mm, fonttitle=\bfseries\color{UTPOrange!85!black},title=\utptitleintuition}
\newtcolorbox{firstread}{enhanced,breakable,colback=white,colframe=UTPBlue!42, boxrule=0.5pt,arc=2mm,left=2.5mm,right=2.5mm,top=1.5mm,bottom=1.5mm, fonttitle=\bfseries\color{UTPBlue}, title=\utptitlefirstread}
\newtcolorbox{debugtask}{enhanced,breakable,colback=UTPPaleRed!58,colframe=UTPRed, boxrule=0.85pt,arc=2mm,left=2.6mm,right=2.6mm,top=1.9mm,bottom=1.9mm, fonttitle=\bfseries\color{white},colbacktitle=UTPRed,title=\utptitledebug}
\newtcolorbox{analogy}{enhanced,breakable,colback=white,colframe=UTPOrange!38, boxrule=0.5pt,arc=2mm,left=2.5mm,right=2.5mm,top=1.5mm,bottom=1.5mm, fonttitle=\bfseries\color{UTPOrange!85!black},title=Analogy}
\newtcolorbox{coursechoice}{enhanced,breakable,colback=white,colframe=UTPNavy, boxrule=0.8pt,arc=2mm,left=2.8mm,right=2.8mm,top=2mm,bottom=2mm, fonttitle=\bfseries\color{white},colbacktitle=UTPNavy,title=Course choice}
\newtcolorbox{decisionmemo}{enhanced,breakable,colback=UTPPaleOrange!55,colframe=UTPOrange, boxrule=0.85pt,arc=2mm,left=2.6mm,right=2.6mm,top=1.9mm,bottom=1.9mm, fonttitle=\bfseries\color{white},colbacktitle=UTPOrange,title=Decision memo}
\newtcolorbox{libnote}{enhanced,breakable,colback=white,colframe=UTPMuted!48, boxrule=0.5pt,arc=2mm,left=2.5mm,right=2.5mm,top=1.5mm,bottom=1.5mm, fonttitle=\bfseries\color{UTPNavy},title=Reference note}
\newtcolorbox{miniproject}{enhanced,breakable,colback=UTPPaleTeal!52,colframe=UTPGreen, boxrule=0.9pt,arc=2mm,left=2.6mm,right=2.6mm,top=1.9mm,bottom=1.9mm, fonttitle=\bfseries\color{white},colbacktitle=UTPGreen,title=Mini-project}
% Boxes that take a title argument
\newtcolorbox{modelcard}[1]{enhanced,breakable,colback=white,colframe=UTPBlue, boxrule=0.9pt,arc=2mm,left=2.5mm,right=2.5mm,top=1.8mm,bottom=1.8mm, fonttitle=\bfseries\color{white},colbacktitle=UTPBlue,title=#1}
\newtcolorbox{optionaltopic}{enhanced,breakable,colback=white,colframe=UTPMuted!58, boxrule=0.5pt,arc=2mm,left=2.5mm,right=2.5mm,top=1.5mm,bottom=1.5mm, fonttitle=\bfseries\color{UTPMuted},title=Beyond the core}
\newtcolorbox{secondpass}{enhanced,breakable,colback=UTPPaleGray!82,colframe=UTPMuted, boxrule=0.75pt,arc=2mm,left=2.6mm,right=2.6mm,top=1.9mm,bottom=1.9mm, fonttitle=\bfseries\color{white},colbacktitle=UTPMuted,title=Second pass}
\newtcolorbox{labproject}[1]{enhanced,breakable,colback=white,colframe=UTPBlue, boxrule=0.9pt,arc=2mm,left=2.6mm,right=2.6mm,top=1.9mm,bottom=1.9mm, fonttitle=\bfseries\color{white},colbacktitle=UTPBlue,title={#1}}
% Frames for a figure or a table
\newtcolorbox{mlfigure}{enhanced,width=\linewidth,colback=white,colframe=UTPBlue!28, boxrule=0.7pt,arc=2mm, left=2.2mm,right=2.2mm,top=2.2mm,bottom=1.8mm, before skip=0.75em,after skip=0.85em, before upper={\centering}, halign=center}
\newtcolorbox{mltable}{enhanced,width=\linewidth,colback=white,colframe=UTPTeal!32, boxrule=0.7pt,arc=2mm, left=2.2mm,right=2.2mm,top=1.8mm,bottom=1.8mm, before skip=0.75em,after skip=0.85em, halign=center}

% --------------------------------------------------------------------
%  9.  SMALL COMMANDS
% --------------------------------------------------------------------
%  \term for a defined term, \cmd for inline code, \chapref and
%  \secref for cross-references, and the mlsteps numbered-step list.
%
%      \term{overfitting}          a term being defined
%      \cmd{git status}            inline code
%      \chapref{ch:trees}          "Chapter 12"
%      \secref{sec:pruning}        "Section 12.3"
%      \chapterguide{prior}{outcomes}   the box that opens a chapter
%      \begin{mlsteps} ... \end{mlsteps}  a numbered Step 1/Step 2 list

\newcommand{\term}[1]{\textbf{\color{UTPNavy}#1}}
\newcommand{\cmd}[1]{\texttt{#1}}
\newcommand{\chapref}[1]{Chapter~\ref{#1}}
\newcommand{\secref}[1]{Section~\ref{#1}}
\newcommand{\chapterguide}[2]{%
  \begin{tcolorbox}[enhanced,breakable,colback=white,colframe=UTPNavy,boxrule=0.8pt,
    arc=2mm,left=2.5mm,right=2.5mm,top=2mm,bottom=2mm,
    title=\utptitleglance,fonttitle=\bfseries\color{white},colbacktitle=UTPNavy]
    \textbf{Before you start:} #1\par
    \textbf{By the end, you can:} #2
  \end{tcolorbox}}
\newenvironment{mlsteps}
  {\begin{enumerate}[label=\textcolor{UTPBlue}{\bfseries Step~\arabic*:},leftmargin=4.5em,itemsep=0.38em]}
  {\end{enumerate}}

% --------------------------------------------------------------------
%  10.  LAB MILESTONE BANNER
% --------------------------------------------------------------------
%  The full-width page break that opens each lab. Called from the lab
%  files as \labsection{number}{title}{label}.
%
%  The "of N" comes from \utplabtotal, so a book states its lab count
%  once in its configuration rather than in every lab file.

\newcommand{\labsection}[3]{%
  \clearpage
  \phantomsection
  \addcontentsline{toc}{lab}{Lab #1 of \utplabtotal\space--- #2}%
  \markboth{Lab #1 of \utplabtotal\space--- #2}{}%
  \begin{tcolorbox}[enhanced,width=\linewidth,colback=UTPNavy,colframe=UTPNavy,
    boxrule=0pt,arc=2mm,left=5mm,right=5mm,top=5mm,bottom=5mm,
    before skip=0pt,after skip=1.2em]
    {\large\bfseries\color{UTPTeal!35!white} PRACTICAL MILESTONE\par}
    \vspace{0.25em}
    {\Huge\bfseries\color{white} Lab #1 of \utplabtotal\par}
    \vspace{0.20em}
    {\Large\bfseries\color{white} #2\par}
  \end{tcolorbox}
  \label{#3}%
}

% --------------------------------------------------------------------
%  11.  STAGE DIVIDER
% --------------------------------------------------------------------
%  A full-page divider used by the longer books to open a part:
%  \substage{3A}{Supervised Workflow}{one-line summary}.
%  Harmless to leave unused.

\newcommand{\substage}[3]{%
  \cleardoublepage
  \thispagestyle{empty}
  \phantomsection
  \addcontentsline{toc}{chapter}{Stage #1 --- #2}
  \begin{center}
    \vspace*{0.20\textheight}
    {\Large\bfseries\color{UTPTeal} Stage #1\par}
    \vspace{0.45em}
    {\Huge\bfseries\color{UTPNavy} #2\par}
    \vspace{0.65em}
    {\color{UTPTeal}\rule{0.46\linewidth}{1.5pt}\par}
    \vspace{1.1em}
    \begin{minipage}{0.76\linewidth}
      \centering\large\color{UTPMuted} #3
    \end{minipage}
  \end{center}
  \cleardoublepage
}

% --------------------------------------------------------------------
%  12.  LISTING STYLES
% --------------------------------------------------------------------
%  Two base styles. A course builds its own on top of them so the
%  colours and frame stay consistent:
%
%      \lstdefinestyle{mlpython}{style=utpcode,language=Python}
%      \lstdefinestyle{mlterminal}{style=utpterminal}

\lstdefinestyle{utpcode}{
  basicstyle=\ttfamily\scriptsize,
  keywordstyle=\bfseries\color{UTPBlue},
  commentstyle=\itshape\color{UTPMuted},
  stringstyle=\color{UTPGreen!70!black},
  backgroundcolor=\color{UTPPaleGray!72},
  frame=single,
  rulecolor=\color{UTPBlue!28},
  framesep=5pt,
  breaklines=true,
  breakatwhitespace=false,
  showstringspaces=false,
  keepspaces=true,
  columns=fullflexible,
  upquote=true,
  tabsize=4,
  aboveskip=0.7em,
  belowskip=0.8em
}
\lstdefinestyle{utpterminal}{
  basicstyle=\ttfamily\scriptsize,
  backgroundcolor=\color{UTPPaleGray!72},
  frame=single,
  rulecolor=\color{UTPMuted!35},
  framesep=5pt,
  breaklines=true,
  showstringspaces=false,
  keepspaces=true,
  columns=fullflexible,
  upquote=true,
  aboveskip=0.6em,
  belowskip=0.7em
}

\lstdefinestyle{adscpp}{
  language=C++,
  basicstyle=\ttfamily\scriptsize,
  keywordstyle=\bfseries\color{UTPBlue},
  commentstyle=\itshape\color{UTPMuted},
  stringstyle=\color{UTPGreen!70!black},
  backgroundcolor=\color{UTPPaleGray!72},
  frame=single,
  rulecolor=\color{UTPBlue!28},
  framesep=5pt,
  breaklines=true,
  breakatwhitespace=false,
  showstringspaces=false,
  keepspaces=true,
  columns=fullflexible,
  upquote=true,
  tabsize=4,
  aboveskip=0.7em,
  belowskip=0.8em
}

\lstdefinestyle{mlpython}{
  language=Python,
  basicstyle=\ttfamily\scriptsize,
  keywordstyle=\bfseries\color{UTPBlue},
  commentstyle=\itshape\color{UTPMuted},
  stringstyle=\color{UTPGreen!70!black},
  backgroundcolor=\color{UTPPaleGray!72},
  frame=single,
  rulecolor=\color{UTPBlue!28},
  framesep=5pt,
  breaklines=true,
  breakatwhitespace=false,
  showstringspaces=false,
  keepspaces=true,
  columns=fullflexible,
  upquote=true,
  tabsize=4,
  aboveskip=0.7em,
  belowskip=0.8em
}

\lstdefinestyle{mlmatlab}{
  language=Matlab,
  basicstyle=\ttfamily\scriptsize,
  keywordstyle=\bfseries\color{UTPBlue},
  commentstyle=\itshape\color{UTPMuted},
  stringstyle=\color{UTPGreen!70!black},
  backgroundcolor=\color{UTPPaleGray!72},
  frame=single,
  rulecolor=\color{UTPBlue!28},
  framesep=5pt,
  breaklines=true,
  breakatwhitespace=false,
  showstringspaces=false,
  keepspaces=true,
  columns=fullflexible,
  upquote=true,
  tabsize=4,
  aboveskip=0.7em,
  belowskip=0.8em
}

\lstdefinestyle{mlterminal}{
  basicstyle=\ttfamily\scriptsize,
  backgroundcolor=\color{UTPPaleGray!72},
  frame=single,
  rulecolor=\color{UTPMuted!35},
  framesep=5pt,
  breaklines=true,
  showstringspaces=false,
  keepspaces=true,
  columns=fullflexible,
  upquote=true,
  aboveskip=0.6em,
  belowskip=0.7em
}

  \lstdefinestyle{dcnterminal}{
  basicstyle=\ttfamily\scriptsize,
  backgroundcolor=\color{UTPPaleGray!72},
  frame=single,rulecolor=\color{UTPMuted!35},
  framesep=5pt,
  breaklines=true,
  showstringspaces=false,
  keepspaces=true,
  columns=fullflexible,
  upquote=true,
  aboveskip=0.6em,
  belowskip=0.7em
}

\lstdefinestyle{dsr}{
  language=R,
  basicstyle=\ttfamily\scriptsize,
  keywordstyle=\bfseries\color{UTPBlue},
  commentstyle=\itshape\color{UTPMuted},
  stringstyle=\color{UTPGreen!70!black},
  backgroundcolor=\color{UTPPaleGray!72},
  frame=single,
  rulecolor=\color{UTPBlue!28},
  framesep=5pt,
  breaklines=true,
  breakatwhitespace=false,
  showstringspaces=false,
  keepspaces=true,
  columns=fullflexible,
  upquote=true,
  tabsize=2,
  aboveskip=0.7em,
  belowskip=0.8em
}

\lstdefinestyle{dsterminal}{
  basicstyle=\ttfamily\scriptsize,
  backgroundcolor=\color{UTPPaleGray!72},
  frame=single,
  rulecolor=\color{UTPMuted!35},
  framesep=5pt,
  breaklines=true,
  showstringspaces=false,
  keepspaces=true,
  columns=fullflexible,
  upquote=true,
  aboveskip=0.6em,
  belowskip=0.7em
}

\lstdefinestyle{erpterminal}{
  basicstyle=\ttfamily\scriptsize,
  backgroundcolor=\color{UTPPaleGray!72},
  frame=single,
  rulecolor=\color{UTPBlue!28},
  framesep=5pt,
  breaklines=true,
  showstringspaces=false,
  keepspaces=true,
  columns=fullflexible,
  upquote=true,
  tabsize=4,
  aboveskip=0.7em,
  belowskip=0.8em
}

\lstdefinestyle{mlpython}{
  language=Python,
  basicstyle=\ttfamily\scriptsize,
  keywordstyle=\bfseries\color{UTPBlue},
  commentstyle=\itshape\color{UTPMuted},
  stringstyle=\color{UTPGreen!70!black},
  backgroundcolor=\color{UTPPaleGray!72},
  frame=single,
  rulecolor=\color{UTPBlue!28},
  framesep=5pt,
  breaklines=true,
  breakatwhitespace=false,
  showstringspaces=false,
  keepspaces=true,
  columns=fullflexible,
  upquote=true,
  tabsize=4,
  aboveskip=0.7em,
  belowskip=0.8em
}
\lstdefinestyle{mlterminal}{
  basicstyle=\ttfamily\scriptsize,
  backgroundcolor=\color{UTPPaleGray!72},
  frame=single,
  rulecolor=\color{UTPMuted!35},
  framesep=5pt,
  breaklines=true,
  showstringspaces=false,
  keepspaces=true,
  columns=fullflexible,
  upquote=true,
  aboveskip=0.6em,
  belowskip=0.7em
}

\lstdefinestyle{mlpython}{
  language=Python,
  basicstyle=\ttfamily\scriptsize,
  keywordstyle=\bfseries\color{UTPBlue},
  commentstyle=\itshape\color{UTPMuted},
  stringstyle=\color{UTPGreen!70!black},
  backgroundcolor=\color{UTPPaleGray!72},
  frame=single,
  rulecolor=\color{UTPBlue!28},
  framesep=5pt,
  breaklines=true,
  breakatwhitespace=false,
  showstringspaces=false,
  keepspaces=true,
  columns=fullflexible,
  upquote=true,
  tabsize=4,
  aboveskip=0.7em,
  belowskip=0.8em
}
\lstdefinestyle{mlterminal}{
  basicstyle=\ttfamily\scriptsize,
  backgroundcolor=\color{UTPPaleGray!72},
  frame=single,
  rulecolor=\color{UTPMuted!35},
  framesep=5pt,
  breaklines=true,
  showstringspaces=false,
  keepspaces=true,
  columns=fullflexible,
  upquote=true,
  aboveskip=0.6em,
  belowskip=0.7em
}
\lstdefinestyle{osc}{
  language=C,
  basicstyle=\ttfamily\scriptsize,
  keywordstyle=\bfseries\color{UTPBlue},
  commentstyle=\itshape\color{UTPMuted},
  stringstyle=\color{UTPGreen!70!black},
  backgroundcolor=\color{UTPPaleGray!72},
  frame=single,
  rulecolor=\color{UTPBlue!28},
  framesep=5pt,
  breaklines=true,
  showstringspaces=false,
  keepspaces=true,
  columns=fullflexible,
  upquote=true,
  tabsize=4,
  aboveskip=0.7em,
  belowskip=0.8em
}
\lstdefinestyle{osbash}{
  language=bash,
  basicstyle=\ttfamily\scriptsize,
  keywordstyle=\bfseries\color{UTPBlue},
  commentstyle=\itshape\color{UTPMuted},
  stringstyle=\color{UTPGreen!70!black},
  backgroundcolor=\color{UTPPaleGray!72},
  frame=single,
  rulecolor=\color{UTPMuted!35},
  framesep=5pt,
  breaklines=true,
  showstringspaces=false,
  keepspaces=true,
  columns=fullflexible,
  upquote=true,
  aboveskip=0.6em,
  belowskip=0.7em
}


